% Môn: Vật lý
% Lớp: 10
% Chương: Chương III. Động lực học
% Bài: L10C3 Bài 19. Lực cản và lực nâng
% Số câu: 162
% App đọc file này trực tiếp — sửa rồi Commit trên GitHub, bấm Đồng bộ GitHub .tex

% L10C3 Bài 19. Lực cản và lực nâng

% ===== Câu 1 =====
% ID: AIQ_L10C3_FULL_0973
% Mức: NB
\dangbt{Nhận biết đặc điểm, phương và chiều của lực cản chất lưu}
\begin{ex}
% ID: AIQ_L10C3_FULL_0973
% Mức: NB
Một vật chuyển động sang phải trong không khí. Lực cản của không khí có chiều nào?
\choice
{\True Sang trái}
{Sang phải}
{Thẳng đứng lên}
{Thẳng đứng xuống}
\loigiai{
Lực cản của chất lưu luôn ngược chiều chuyển động tương đối của vật đối với chất lưu. Chọn A.
}
\end{ex}

% ===== Câu 2 =====
% ID: AIQ_L10C3_FULL_0974
% Mức: NB
\begin{ex}
% ID: AIQ_L10C3_FULL_0974
% Mức: NB
Một vật chuyển động sang phải trong không khí. Lực cản của không khí có chiều nào?
\choice
{\True Sang trái}
{Sang phải}
{Thẳng đứng lên}
{Thẳng đứng xuống}
\loigiai{
Lực cản của chất lưu luôn ngược chiều chuyển động tương đối của vật đối với chất lưu. Chọn A.
}
\end{ex}

% ===== Câu 3 =====
% ID: AIQ_L10C3_FULL_0975
% Mức: NB
\begin{ex}
% ID: AIQ_L10C3_FULL_0975
% Mức: NB
Một vật chuyển động sang phải trong không khí. Lực cản của không khí có chiều nào?
\choice
{\True Sang trái}
{Sang phải}
{Thẳng đứng lên}
{Thẳng đứng xuống}
\loigiai{
Lực cản của chất lưu luôn ngược chiều chuyển động tương đối của vật đối với chất lưu. Chọn A.
}
\end{ex}

% ===== Câu 4 =====
% ID: AIQ_L10C3_FULL_0976
% Mức: TH
\begin{ex}
% ID: AIQ_L10C3_FULL_0976
% Mức: TH
Một vật chuyển động sang phải trong không khí. Lực cản của không khí có chiều nào?
\choice
{\True Sang trái}
{Sang phải}
{Thẳng đứng lên}
{Thẳng đứng xuống}
\loigiai{
Lực cản của chất lưu luôn ngược chiều chuyển động tương đối của vật đối với chất lưu. Chọn A.
}
\end{ex}

% ===== Câu 5 =====
% ID: AIQ_L10C3_FULL_0977
% Mức: TH
\begin{ex}
% ID: AIQ_L10C3_FULL_0977
% Mức: TH
Một vật chuyển động sang phải trong không khí. Lực cản của không khí có chiều nào?
\choice
{\True Sang trái}
{Sang phải}
{Thẳng đứng lên}
{Thẳng đứng xuống}
\loigiai{
Lực cản của chất lưu luôn ngược chiều chuyển động tương đối của vật đối với chất lưu. Chọn A.
}
\end{ex}

% ===== Câu 6 =====
% ID: AIQ_L10C3_FULL_0978
% Mức: TH
\begin{ex}
% ID: AIQ_L10C3_FULL_0978
% Mức: TH
Một vật chuyển động sang phải trong không khí. Lực cản của không khí có chiều nào?
\choice
{\True Sang trái}
{Sang phải}
{Thẳng đứng lên}
{Thẳng đứng xuống}
\loigiai{
Lực cản của chất lưu luôn ngược chiều chuyển động tương đối của vật đối với chất lưu. Chọn A.
}
\end{ex}

% ===== Câu 7 =====
% ID: AIQ_L10C3_FULL_0979
% Mức: TH
\begin{ex}
% ID: AIQ_L10C3_FULL_0979
% Mức: TH
Một vật chuyển động sang phải trong không khí. Lực cản của không khí có chiều nào?
\choice
{\True Sang trái}
{Sang phải}
{Thẳng đứng lên}
{Thẳng đứng xuống}
\loigiai{
Lực cản của chất lưu luôn ngược chiều chuyển động tương đối của vật đối với chất lưu. Chọn A.
}
\end{ex}

% ===== Câu 8 =====
% ID: AIQ_L10C3_FULL_0980
% Mức: VD
\begin{ex}
% ID: AIQ_L10C3_FULL_0980
% Mức: VD
Một vật chuyển động sang phải trong không khí. Lực cản của không khí có chiều nào?
\choice
{\True Sang trái}
{Sang phải}
{Thẳng đứng lên}
{Thẳng đứng xuống}
\loigiai{
Lực cản của chất lưu luôn ngược chiều chuyển động tương đối của vật đối với chất lưu. Chọn A.
}
\end{ex}

% ===== Câu 9 =====
% ID: AIQ_L10C3_FULL_0981
% Mức: VD
\begin{ex}
% ID: AIQ_L10C3_FULL_0981
% Mức: VD
Một vật chuyển động sang phải trong không khí. Lực cản của không khí có chiều nào?
\choice
{\True Sang trái}
{Sang phải}
{Thẳng đứng lên}
{Thẳng đứng xuống}
\loigiai{
Lực cản của chất lưu luôn ngược chiều chuyển động tương đối của vật đối với chất lưu. Chọn A.
}
\end{ex}

% ===== Câu 10 =====
% ID: AIQ_L10C3_FULL_0982
% Mức: VDC
\begin{ex}
% ID: AIQ_L10C3_FULL_0982
% Mức: VDC
Một vật chuyển động sang phải trong không khí. Lực cản của không khí có chiều nào?
\choice
{\True Sang trái}
{Sang phải}
{Thẳng đứng lên}
{Thẳng đứng xuống}
\loigiai{
Lực cản của chất lưu luôn ngược chiều chuyển động tương đối của vật đối với chất lưu. Chọn A.
}
\end{ex}

% ===== Câu 11 =====
% ID: AIQ_L10C3_FULL_0983
% Mức: NB
\begin{ex}
% ID: AIQ_L10C3_FULL_0983
% Mức: NB
Một vật chuyển động sang phải trong không khí. Lực cản của không khí có chiều nào?
\choiceTF
{\True Lực cản phụ thuộc chuyển động tương đối giữa vật và chất lưu.}
{Lực cản luôn cùng chiều chuyển động của vật.}
{\True Khi hợp lực bằng 0, vật có gia tốc bằng 0.}
{Có thể bỏ qua đơn vị của đại lượng trong kết quả vật lí.}
\loigiai{
\begin{itemchoice}
\item (Đúng) Lực cản phụ thuộc chuyển động tương đối giữa vật và chất lưu. (Vì): Lực cản của chất lưu luôn ngược chiều chuyển động tương đối của vật đối với chất lưu. b) Sai vì giá trị hoặc chiều nêu ra không phù hợp. c) Đúng theo định luật II Newton. d) Sai vì kết quả vật lí phải kèm đơn vị thích hợp
\item (Sai) Lực cản luôn cùng chiều chuyển động của vật.
\item (Đúng) Khi hợp lực bằng 0, vật có gia tốc bằng 0.
\item (Sai) Có thể bỏ qua đơn vị của đại lượng trong kết quả vật lí.
\end{itemchoice}
}
\end{ex}

% ===== Câu 12 =====
% ID: AIQ_L10C3_FULL_0984
% Mức: NB
\begin{ex}
% ID: AIQ_L10C3_FULL_0984
% Mức: NB
Một vật chuyển động sang phải trong không khí. Lực cản của không khí có chiều nào?
\choiceTF
{\True Lực cản phụ thuộc chuyển động tương đối giữa vật và chất lưu.}
{Lực cản luôn cùng chiều chuyển động của vật.}
{\True Khi hợp lực bằng 0, vật có gia tốc bằng 0.}
{Có thể bỏ qua đơn vị của đại lượng trong kết quả vật lí.}
\loigiai{
\begin{itemchoice}
\item (Đúng) Lực cản phụ thuộc chuyển động tương đối giữa vật và chất lưu. (Vì): Lực cản của chất lưu luôn ngược chiều chuyển động tương đối của vật đối với chất lưu. b) Sai vì giá trị hoặc chiều nêu ra không phù hợp. c) Đúng theo định luật II Newton. d) Sai vì kết quả vật lí phải kèm đơn vị thích hợp
\item (Sai) Lực cản luôn cùng chiều chuyển động của vật.
\item (Đúng) Khi hợp lực bằng 0, vật có gia tốc bằng 0.
\item (Sai) Có thể bỏ qua đơn vị của đại lượng trong kết quả vật lí.
\end{itemchoice}
}
\end{ex}

% ===== Câu 13 =====
% ID: AIQ_L10C3_FULL_0985
% Mức: TH
\begin{ex}
% ID: AIQ_L10C3_FULL_0985
% Mức: TH
Một vật chuyển động sang phải trong không khí. Lực cản của không khí có chiều nào?
\choiceTF
{\True Lực cản phụ thuộc chuyển động tương đối giữa vật và chất lưu.}
{Lực cản luôn cùng chiều chuyển động của vật.}
{\True Khi hợp lực bằng 0, vật có gia tốc bằng 0.}
{Có thể bỏ qua đơn vị của đại lượng trong kết quả vật lí.}
\loigiai{
\begin{itemchoice}
\item (Đúng) Lực cản phụ thuộc chuyển động tương đối giữa vật và chất lưu. (Vì): Lực cản của chất lưu luôn ngược chiều chuyển động tương đối của vật đối với chất lưu. b) Sai vì giá trị hoặc chiều nêu ra không phù hợp. c) Đúng theo định luật II Newton. d) Sai vì kết quả vật lí phải kèm đơn vị thích hợp
\item (Sai) Lực cản luôn cùng chiều chuyển động của vật.
\item (Đúng) Khi hợp lực bằng 0, vật có gia tốc bằng 0.
\item (Sai) Có thể bỏ qua đơn vị của đại lượng trong kết quả vật lí.
\end{itemchoice}
}
\end{ex}

% ===== Câu 14 =====
% ID: AIQ_L10C3_FULL_0986
% Mức: TH
\begin{ex}
% ID: AIQ_L10C3_FULL_0986
% Mức: TH
Một vật chuyển động sang phải trong không khí. Lực cản của không khí có chiều nào?
\choiceTF
{\True Lực cản phụ thuộc chuyển động tương đối giữa vật và chất lưu.}
{Lực cản luôn cùng chiều chuyển động của vật.}
{\True Khi hợp lực bằng 0, vật có gia tốc bằng 0.}
{Có thể bỏ qua đơn vị của đại lượng trong kết quả vật lí.}
\loigiai{
\begin{itemchoice}
\item (Đúng) Lực cản phụ thuộc chuyển động tương đối giữa vật và chất lưu. (Vì): Lực cản của chất lưu luôn ngược chiều chuyển động tương đối của vật đối với chất lưu. b) Sai vì giá trị hoặc chiều nêu ra không phù hợp. c) Đúng theo định luật II Newton. d) Sai vì kết quả vật lí phải kèm đơn vị thích hợp
\item (Sai) Lực cản luôn cùng chiều chuyển động của vật.
\item (Đúng) Khi hợp lực bằng 0, vật có gia tốc bằng 0.
\item (Sai) Có thể bỏ qua đơn vị của đại lượng trong kết quả vật lí.
\end{itemchoice}
}
\end{ex}

% ===== Câu 15 =====
% ID: AIQ_L10C3_FULL_0987
% Mức: VD
\begin{ex}
% ID: AIQ_L10C3_FULL_0987
% Mức: VD
Một vật chuyển động sang phải trong không khí. Lực cản của không khí có chiều nào?
\choiceTF
{\True Lực cản phụ thuộc chuyển động tương đối giữa vật và chất lưu.}
{Lực cản luôn cùng chiều chuyển động của vật.}
{\True Khi hợp lực bằng 0, vật có gia tốc bằng 0.}
{Có thể bỏ qua đơn vị của đại lượng trong kết quả vật lí.}
\loigiai{
\begin{itemchoice}
\item (Đúng) Lực cản phụ thuộc chuyển động tương đối giữa vật và chất lưu. (Vì): Lực cản của chất lưu luôn ngược chiều chuyển động tương đối của vật đối với chất lưu. b) Sai vì giá trị hoặc chiều nêu ra không phù hợp. c) Đúng theo định luật II Newton. d) Sai vì kết quả vật lí phải kèm đơn vị thích hợp
\item (Sai) Lực cản luôn cùng chiều chuyển động của vật.
\item (Đúng) Khi hợp lực bằng 0, vật có gia tốc bằng 0.
\item (Sai) Có thể bỏ qua đơn vị của đại lượng trong kết quả vật lí.
\end{itemchoice}
}
\end{ex}

% ===== Câu 16 =====
% ID: AIQ_L10C3_FULL_0988
% Mức: TH
\begin{ex}
% ID: AIQ_L10C3_FULL_0988
% Mức: TH
Một vật chuyển động sang phải trong không khí. Lực cản của không khí có chiều nào?
\shortans{Sang}
\loigiai{
Lực cản của chất lưu luôn ngược chiều chuyển động tương đối của vật đối với chất lưu. Đáp án trả lời ngắn không ghi đơn vị.
}
\end{ex}

% ===== Câu 17 =====
% ID: AIQ_L10C3_FULL_0989
% Mức: TH
\begin{ex}
% ID: AIQ_L10C3_FULL_0989
% Mức: TH
Một vật chuyển động sang phải trong không khí. Lực cản của không khí có chiều nào?
\shortans{Sang}
\loigiai{
Lực cản của chất lưu luôn ngược chiều chuyển động tương đối của vật đối với chất lưu. Đáp án trả lời ngắn không ghi đơn vị.
}
\end{ex}

% ===== Câu 18 =====
% ID: AIQ_L10C3_FULL_0990
% Mức: VD
\begin{ex}
% ID: AIQ_L10C3_FULL_0990
% Mức: VD
Một vật chuyển động sang phải trong không khí. Lực cản của không khí có chiều nào?
\shortans{Sang}
\loigiai{
Lực cản của chất lưu luôn ngược chiều chuyển động tương đối của vật đối với chất lưu. Đáp án trả lời ngắn không ghi đơn vị.
}
\end{ex}

% ===== Câu 19 =====
% ID: AIQ_L10C3_FULL_0991
% Mức: VD
\begin{ex}
% ID: AIQ_L10C3_FULL_0991
% Mức: VD
Một vật chuyển động sang phải trong không khí. Lực cản của không khí có chiều nào?
\shortans{Sang}
\loigiai{
Lực cản của chất lưu luôn ngược chiều chuyển động tương đối của vật đối với chất lưu. Đáp án trả lời ngắn không ghi đơn vị.
}
\end{ex}

% ===== Câu 20 =====
% ID: AIQ_L10C3_FULL_0992
% Mức: VD
\begin{ex}
% ID: AIQ_L10C3_FULL_0992
% Mức: VD
Một vật chuyển động sang phải trong không khí. Lực cản của không khí có chiều nào?
\shortans{Sang}
\loigiai{
Lực cản của chất lưu luôn ngược chiều chuyển động tương đối của vật đối với chất lưu. Đáp án trả lời ngắn không ghi đơn vị.
}
\end{ex}

% ===== Câu 21 =====
% ID: AIQ_L10C3_FULL_0993
% Mức: VDC
\begin{ex}
% ID: AIQ_L10C3_FULL_0993
% Mức: VDC
Một vật chuyển động sang phải trong không khí. Lực cản của không khí có chiều nào?
\shortans{Sang}
\loigiai{
Lực cản của chất lưu luôn ngược chiều chuyển động tương đối của vật đối với chất lưu. Đáp án trả lời ngắn không ghi đơn vị.
}
\end{ex}

% ===== Câu 22 =====
% ID: AIQ_L10C3_FULL_0994
% Mức: TH
\begin{bt}
% ID: AIQ_L10C3_FULL_0994
% Mức: TH
Một vật chuyển động sang phải trong không khí. Lực cản của không khí có chiều nào?
\loigiai{
Phân tích dữ kiện, chọn định luật phù hợp. Lực cản của chất lưu luôn ngược chiều chuyển động tương đối của vật đối với chất lưu.
}
\end{bt}

% ===== Câu 23 =====
% ID: AIQ_L10C3_FULL_0995
% Mức: TH
\begin{bt}
% ID: AIQ_L10C3_FULL_0995
% Mức: TH
Một vật chuyển động sang phải trong không khí. Lực cản của không khí có chiều nào?
\loigiai{
Phân tích dữ kiện, chọn định luật phù hợp. Lực cản của chất lưu luôn ngược chiều chuyển động tương đối của vật đối với chất lưu.
}
\end{bt}

% ===== Câu 24 =====
% ID: AIQ_L10C3_FULL_0996
% Mức: VD
\begin{bt}
% ID: AIQ_L10C3_FULL_0996
% Mức: VD
Một vật chuyển động sang phải trong không khí. Lực cản của không khí có chiều nào?
\loigiai{
Phân tích dữ kiện, chọn định luật phù hợp. Lực cản của chất lưu luôn ngược chiều chuyển động tương đối của vật đối với chất lưu.
}
\end{bt}

% ===== Câu 25 =====
% ID: AIQ_L10C3_FULL_0997
% Mức: VD
\begin{bt}
% ID: AIQ_L10C3_FULL_0997
% Mức: VD
Một vật chuyển động sang phải trong không khí. Lực cản của không khí có chiều nào?
\loigiai{
Phân tích dữ kiện, chọn định luật phù hợp. Lực cản của chất lưu luôn ngược chiều chuyển động tương đối của vật đối với chất lưu.
}
\end{bt}

% ===== Câu 26 =====
% ID: AIQ_L10C3_FULL_0998
% Mức: VD
\begin{bt}
% ID: AIQ_L10C3_FULL_0998
% Mức: VD
Một vật chuyển động sang phải trong không khí. Lực cản của không khí có chiều nào?
\loigiai{
Phân tích dữ kiện, chọn định luật phù hợp. Lực cản của chất lưu luôn ngược chiều chuyển động tương đối của vật đối với chất lưu.
}
\end{bt}

% ===== Câu 27 =====
% ID: AIQ_L10C3_FULL_0999
% Mức: VDC
\begin{bt}
% ID: AIQ_L10C3_FULL_0999
% Mức: VDC
Một vật chuyển động sang phải trong không khí. Lực cản của không khí có chiều nào?
\loigiai{
Phân tích dữ kiện, chọn định luật phù hợp. Lực cản của chất lưu luôn ngược chiều chuyển động tương đối của vật đối với chất lưu.
}
\end{bt}

% ===== Câu 28 =====
% ID: AIQ_L10C3_FULL_1000
% Mức: NB
\dangbt{Tính lực cản của chất lưu theo tốc độ chuyển động}
\begin{ex}
% ID: AIQ_L10C3_FULL_1000
% Mức: NB
Vật chuyển động trong không khí với tốc độ $5\,\text{m}$ /s. Biết lực cản $F_c=\frac{1}{2}\rho C A v^2$, $\rho=1,2$ kg/m^3, $C=0,5$, $A=0,2$ m^2. Tính lực cản.
\choice
{\True 1,5 $N$}
{1,8 $N$}
{1,2 $N$}
{3,5 $N$}
\loigiai{
Thay số: $F_c=\frac12\cdot1,2\cdot0,5\cdot0,2\cdot5^2=1,5$ N. Chọn A.
}
\end{ex}

% ===== Câu 29 =====
% ID: AIQ_L10C3_FULL_1001
% Mức: NB
\begin{ex}
% ID: AIQ_L10C3_FULL_1001
% Mức: NB
Vật chuyển động trong không khí với tốc độ $6\,\text{m}$ /s. Biết lực cản $F_c=\frac{1}{2}\rho C A v^2$, $\rho=1,2$ kg/m^3, $C=0,6$, $A=0,25$ m^2. Tính lực cản.
\choice
{3,89 $N$}
{2,59 $N$}
{5,24 $N$}
{\True 3,24 $N$}
\loigiai{
Thay số: $F_c=\frac12\cdot1,2\cdot0,6\cdot0,25\cdot6^2=3,24$ N. Chọn D.
}
\end{ex}

% ===== Câu 30 =====
% ID: AIQ_L10C3_FULL_1002
% Mức: NB
\begin{ex}
% ID: AIQ_L10C3_FULL_1002
% Mức: NB
Vật chuyển động trong không khí với tốc độ $7\,\text{m}$ /s. Biết lực cản $F_c=\frac{1}{2}\rho C A v^2$, $\rho=1,2$ kg/m^3, $C=0,4$, $A=0,3$ m^2. Tính lực cản.
\choice
{2,82 $N$}
{5,53 $N$}
{\True 3,53 $N$}
{4,24 $N$}
\loigiai{
Thay số: $F_c=\frac12\cdot1,2\cdot0,4\cdot0,3\cdot7^2=3,53$ N. Chọn C.
}
\end{ex}

% ===== Câu 31 =====
% ID: AIQ_L10C3_FULL_1003
% Mức: TH
\begin{ex}
% ID: AIQ_L10C3_FULL_1003
% Mức: TH
Vật chuyển động trong không khí với tốc độ $8\,\text{m}$ /s. Biết lực cản $F_c=\frac{1}{2}\rho C A v^2$, $\rho=1,2$ kg/m^3, $C=0,5$, $A=0,35$ m^2. Tính lực cản.
\choice
{8,72 $N$}
{\True 6,72 $N$}
{8,06 $N$}
{5,38 $N$}
\loigiai{
Thay số: $F_c=\frac12\cdot1,2\cdot0,5\cdot0,35\cdot8^2=6,72$ N. Chọn B.
}
\end{ex}

% ===== Câu 32 =====
% ID: AIQ_L10C3_FULL_1004
% Mức: TH
\begin{ex}
% ID: AIQ_L10C3_FULL_1004
% Mức: TH
Vật chuyển động trong không khí với tốc độ $9\,\text{m}$ /s. Biết lực cản $F_c=\frac{1}{2}\rho C A v^2$, $\rho=1,2$ kg/m^3, $C=0,6$, $A=0,4$ m^2. Tính lực cản.
\choice
{\True 11,66 $N$}
{13,99 $N$}
{9,33 $N$}
{13,66 $N$}
\loigiai{
Thay số: $F_c=\frac12\cdot1,2\cdot0,6\cdot0,4\cdot9^2=11,66$ N. Chọn A.
}
\end{ex}

% ===== Câu 33 =====
% ID: AIQ_L10C3_FULL_1005
% Mức: TH
\begin{ex}
% ID: AIQ_L10C3_FULL_1005
% Mức: TH
Vật chuyển động trong không khí với tốc độ $10\,\text{m}$ /s. Biết lực cản $F_c=\frac{1}{2}\rho C A v^2$, $\rho=1,2$ kg/m^3, $C=0,4$, $A=0,2$ m^2. Tính lực cản.
\choice
{5,76 $N$}
{3,84 $N$}
{6,8 $N$}
{\True 4,8 $N$}
\loigiai{
Thay số: $F_c=\frac12\cdot1,2\cdot0,4\cdot0,2\cdot10^2=4,8$ N. Chọn D.
}
\end{ex}

% ===== Câu 34 =====
% ID: AIQ_L10C3_FULL_1006
% Mức: TH
\begin{ex}
% ID: AIQ_L10C3_FULL_1006
% Mức: TH
Vật chuyển động trong không khí với tốc độ $11\,\text{m}$ /s. Biết lực cản $F_c=\frac{1}{2}\rho C A v^2$, $\rho=1,2$ kg/m^3, $C=0,5$, $A=0,25$ m^2. Tính lực cản.
\choice
{7,26 $N$}
{11,07 $N$}
{\True 9,07 $N$}
{10,88 $N$}
\loigiai{
Thay số: $F_c=\frac12\cdot1,2\cdot0,5\cdot0,25\cdot11^2=9,07$ N. Chọn C.
}
\end{ex}

% ===== Câu 35 =====
% ID: AIQ_L10C3_FULL_1007
% Mức: VD
\begin{ex}
% ID: AIQ_L10C3_FULL_1007
% Mức: VD
Vật chuyển động trong không khí với tốc độ $12\,\text{m}$ /s. Biết lực cản $F_c=\frac{1}{2}\rho C A v^2$, $\rho=1,2$ kg/m^3, $C=0,6$, $A=0,3$ m^2. Tính lực cản.
\choice
{17,55 $N$}
{\True 15,55 $N$}
{18,66 $N$}
{12,44 $N$}
\loigiai{
Thay số: $F_c=\frac12\cdot1,2\cdot0,6\cdot0,3\cdot12^2=15,55$ N. Chọn B.
}
\end{ex}

% ===== Câu 36 =====
% ID: AIQ_L10C3_FULL_1008
% Mức: VD
\begin{ex}
% ID: AIQ_L10C3_FULL_1008
% Mức: VD
Vật chuyển động trong không khí với tốc độ $5\,\text{m}$ /s. Biết lực cản $F_c=\frac{1}{2}\rho C A v^2$, $\rho=1,2$ kg/m^3, $C=0,4$, $A=0,35$ m^2. Tính lực cản.
\choice
{\True 2,1 $N$}
{2,52 $N$}
{1,68 $N$}
{4,1 $N$}
\loigiai{
Thay số: $F_c=\frac12\cdot1,2\cdot0,4\cdot0,35\cdot5^2=2,1$ N. Chọn A.
}
\end{ex}

% ===== Câu 37 =====
% ID: AIQ_L10C3_FULL_1009
% Mức: VDC
\begin{ex}
% ID: AIQ_L10C3_FULL_1009
% Mức: VDC
Vật chuyển động trong không khí với tốc độ $6\,\text{m}$ /s. Biết lực cản $F_c=\frac{1}{2}\rho C A v^2$, $\rho=1,2$ kg/m^3, $C=0,5$, $A=0,4$ m^2. Tính lực cản.
\choice
{5,18 $N$}
{3,46 $N$}
{6,32 $N$}
{\True 4,32 $N$}
\loigiai{
Thay số: $F_c=\frac12\cdot1,2\cdot0,5\cdot0,4\cdot6^2=4,32$ N. Chọn D.
}
\end{ex}

% ===== Câu 38 =====
% ID: AIQ_L10C3_FULL_1010
% Mức: NB
\begin{ex}
% ID: AIQ_L10C3_FULL_1010
% Mức: NB
Vật chuyển động trong không khí với tốc độ $7\,\text{m}$ /s. Biết lực cản $F_c=\frac{1}{2}\rho C A v^2$, $\rho=1,2$ kg/m^3, $C=0,6$, $A=0,2$ m^2. Tính lực cản.
\choiceTF
{\True Kết quả đúng là $3,53\,\text{N}$.}
{Kết quả bằng $7,06\,\text{N}$.}
{\True Khi hợp lực bằng 0, vật có gia tốc bằng 0.}
{Có thể bỏ qua đơn vị của đại lượng trong kết quả vật lí.}
\loigiai{
\begin{itemchoice}
\item (Đúng) Kết quả đúng là $3,53\,\text{N}$. (Vì): Thay số: $F_c=\frac12\cdot1,2\cdot0,6\cdot0,2\cdot7^2=3,53$ N. b) Sai vì giá trị hoặc chiều nêu ra không phù hợp. c) Đúng theo định luật II Newton. d) Sai vì kết quả vật lí phải kèm đơn vị thích hợp
\item (Sai) Kết quả bằng $7,06\,\text{N}$.
\item (Đúng) Khi hợp lực bằng 0, vật có gia tốc bằng 0.
\item (Sai) Có thể bỏ qua đơn vị của đại lượng trong kết quả vật lí.
\end{itemchoice}
}
\end{ex}

% ===== Câu 39 =====
% ID: AIQ_L10C3_FULL_1011
% Mức: NB
\begin{ex}
% ID: AIQ_L10C3_FULL_1011
% Mức: NB
Vật chuyển động trong không khí với tốc độ $8\,\text{m}$ /s. Biết lực cản $F_c=\frac{1}{2}\rho C A v^2$, $\rho=1,2$ kg/m^3, $C=0,4$, $A=0,25$ m^2. Tính lực cản.
\choiceTF
{\True Kết quả đúng là $3,84\,\text{N}$.}
{Kết quả bằng $7,68\,\text{N}$.}
{\True Khi hợp lực bằng 0, vật có gia tốc bằng 0.}
{Có thể bỏ qua đơn vị của đại lượng trong kết quả vật lí.}
\loigiai{
\begin{itemchoice}
\item (Đúng) Kết quả đúng là $3,84\,\text{N}$. (Vì): Thay số: $F_c=\frac12\cdot1,2\cdot0,4\cdot0,25\cdot8^2=3,84$ N. b) Sai vì giá trị hoặc chiều nêu ra không phù hợp. c) Đúng theo định luật II Newton. d) Sai vì kết quả vật lí phải kèm đơn vị thích hợp
\item (Sai) Kết quả bằng $7,68\,\text{N}$.
\item (Đúng) Khi hợp lực bằng 0, vật có gia tốc bằng 0.
\item (Sai) Có thể bỏ qua đơn vị của đại lượng trong kết quả vật lí.
\end{itemchoice}
}
\end{ex}

% ===== Câu 40 =====
% ID: AIQ_L10C3_FULL_1012
% Mức: TH
\begin{ex}
% ID: AIQ_L10C3_FULL_1012
% Mức: TH
Vật chuyển động trong không khí với tốc độ $9\,\text{m}$ /s. Biết lực cản $F_c=\frac{1}{2}\rho C A v^2$, $\rho=1,2$ kg/m^3, $C=0,5$, $A=0,3$ m^2. Tính lực cản.
\choiceTF
{\True Kết quả đúng là $7,29\,\text{N}$.}
{Kết quả bằng $14,58\,\text{N}$.}
{\True Khi hợp lực bằng 0, vật có gia tốc bằng 0.}
{Có thể bỏ qua đơn vị của đại lượng trong kết quả vật lí.}
\loigiai{
\begin{itemchoice}
\item (Đúng) Kết quả đúng là $7,29\,\text{N}$. (Vì): Thay số: $F_c=\frac12\cdot1,2\cdot0,5\cdot0,3\cdot9^2=7,29$ N. b) Sai vì giá trị hoặc chiều nêu ra không phù hợp. c) Đúng theo định luật II Newton. d) Sai vì kết quả vật lí phải kèm đơn vị thích hợp
\item (Sai) Kết quả bằng $14,58\,\text{N}$.
\item (Đúng) Khi hợp lực bằng 0, vật có gia tốc bằng 0.
\item (Sai) Có thể bỏ qua đơn vị của đại lượng trong kết quả vật lí.
\end{itemchoice}
}
\end{ex}

% ===== Câu 41 =====
% ID: AIQ_L10C3_FULL_1013
% Mức: TH
\begin{ex}
% ID: AIQ_L10C3_FULL_1013
% Mức: TH
Vật chuyển động trong không khí với tốc độ $10\,\text{m}$ /s. Biết lực cản $F_c=\frac{1}{2}\rho C A v^2$, $\rho=1,2$ kg/m^3, $C=0,6$, $A=0,35$ m^2. Tính lực cản.
\choiceTF
{\True Kết quả đúng là $12,6\,\text{N}$.}
{Kết quả bằng $25,2\,\text{N}$.}
{\True Khi hợp lực bằng 0, vật có gia tốc bằng 0.}
{Có thể bỏ qua đơn vị của đại lượng trong kết quả vật lí.}
\loigiai{
\begin{itemchoice}
\item (Đúng) Kết quả đúng là $12,6\,\text{N}$. (Vì): Thay số: $F_c=\frac12\cdot1,2\cdot0,6\cdot0,35\cdot10^2=12,6$ N. b) Sai vì giá trị hoặc chiều nêu ra không phù hợp. c) Đúng theo định luật II Newton. d) Sai vì kết quả vật lí phải kèm đơn vị thích hợp
\item (Sai) Kết quả bằng $25,2\,\text{N}$.
\item (Đúng) Khi hợp lực bằng 0, vật có gia tốc bằng 0.
\item (Sai) Có thể bỏ qua đơn vị của đại lượng trong kết quả vật lí.
\end{itemchoice}
}
\end{ex}

% ===== Câu 42 =====
% ID: AIQ_L10C3_FULL_1014
% Mức: VD
\begin{ex}
% ID: AIQ_L10C3_FULL_1014
% Mức: VD
Vật chuyển động trong không khí với tốc độ $11\,\text{m}$ /s. Biết lực cản $F_c=\frac{1}{2}\rho C A v^2$, $\rho=1,2$ kg/m^3, $C=0,4$, $A=0,4$ m^2. Tính lực cản.
\choiceTF
{\True Kết quả đúng là $11,62\,\text{N}$.}
{Kết quả bằng $23,24\,\text{N}$.}
{\True Khi hợp lực bằng 0, vật có gia tốc bằng 0.}
{Có thể bỏ qua đơn vị của đại lượng trong kết quả vật lí.}
\loigiai{
\begin{itemchoice}
\item (Đúng) Kết quả đúng là $11,62\,\text{N}$. (Vì): Thay số: $F_c=\frac12\cdot1,2\cdot0,4\cdot0,4\cdot11^2=11,62$ N. b) Sai vì giá trị hoặc chiều nêu ra không phù hợp. c) Đúng theo định luật II Newton. d) Sai vì kết quả vật lí phải kèm đơn vị thích hợp
\item (Sai) Kết quả bằng $23,24\,\text{N}$.
\item (Đúng) Khi hợp lực bằng 0, vật có gia tốc bằng 0.
\item (Sai) Có thể bỏ qua đơn vị của đại lượng trong kết quả vật lí.
\end{itemchoice}
}
\end{ex}

% ===== Câu 43 =====
% ID: AIQ_L10C3_FULL_1015
% Mức: TH
\begin{ex}
% ID: AIQ_L10C3_FULL_1015
% Mức: TH
Vật chuyển động trong không khí với tốc độ $12\,\text{m}$ /s. Biết lực cản $F_c=\frac{1}{2}\rho C A v^2$, $\rho=1,2$ kg/m^3, $C=0,5$, $A=0,2$ m^2. Tính lực cản.
\shortans{8,64}
\loigiai{
Thay số: $F_c=\frac12\cdot1,2\cdot0,5\cdot0,2\cdot12^2=8,64$ N. Đáp án trả lời ngắn không ghi đơn vị.
}
\end{ex}

% ===== Câu 44 =====
% ID: AIQ_L10C3_FULL_1016
% Mức: TH
\begin{ex}
% ID: AIQ_L10C3_FULL_1016
% Mức: TH
Vật chuyển động trong không khí với tốc độ $5\,\text{m}$ /s. Biết lực cản $F_c=\frac{1}{2}\rho C A v^2$, $\rho=1,2$ kg/m^3, $C=0,6$, $A=0,25$ m^2. Tính lực cản.
\shortans{2,25}
\loigiai{
Thay số: $F_c=\frac12\cdot1,2\cdot0,6\cdot0,25\cdot5^2=2,25$ N. Đáp án trả lời ngắn không ghi đơn vị.
}
\end{ex}

% ===== Câu 45 =====
% ID: AIQ_L10C3_FULL_1017
% Mức: VD
\begin{ex}
% ID: AIQ_L10C3_FULL_1017
% Mức: VD
Vật chuyển động trong không khí với tốc độ $6\,\text{m}$ /s. Biết lực cản $F_c=\frac{1}{2}\rho C A v^2$, $\rho=1,2$ kg/m^3, $C=0,4$, $A=0,3$ m^2. Tính lực cản.
\shortans{2,59}
\loigiai{
Thay số: $F_c=\frac12\cdot1,2\cdot0,4\cdot0,3\cdot6^2=2,59$ N. Đáp án trả lời ngắn không ghi đơn vị.
}
\end{ex}

% ===== Câu 46 =====
% ID: AIQ_L10C3_FULL_1018
% Mức: VD
\begin{ex}
% ID: AIQ_L10C3_FULL_1018
% Mức: VD
Vật chuyển động trong không khí với tốc độ $7\,\text{m}$ /s. Biết lực cản $F_c=\frac{1}{2}\rho C A v^2$, $\rho=1,2$ kg/m^3, $C=0,5$, $A=0,35$ m^2. Tính lực cản.
\shortans{5,15}
\loigiai{
Thay số: $F_c=\frac12\cdot1,2\cdot0,5\cdot0,35\cdot7^2=5,15$ N. Đáp án trả lời ngắn không ghi đơn vị.
}
\end{ex}

% ===== Câu 47 =====
% ID: AIQ_L10C3_FULL_1019
% Mức: VD
\begin{ex}
% ID: AIQ_L10C3_FULL_1019
% Mức: VD
Vật chuyển động trong không khí với tốc độ $8\,\text{m}$ /s. Biết lực cản $F_c=\frac{1}{2}\rho C A v^2$, $\rho=1,2$ kg/m^3, $C=0,6$, $A=0,4$ m^2. Tính lực cản.
\shortans{9,22}
\loigiai{
Thay số: $F_c=\frac12\cdot1,2\cdot0,6\cdot0,4\cdot8^2=9,22$ N. Đáp án trả lời ngắn không ghi đơn vị.
}
\end{ex}

% ===== Câu 48 =====
% ID: AIQ_L10C3_FULL_1020
% Mức: VDC
\begin{ex}
% ID: AIQ_L10C3_FULL_1020
% Mức: VDC
Vật chuyển động trong không khí với tốc độ $9\,\text{m}$ /s. Biết lực cản $F_c=\frac{1}{2}\rho C A v^2$, $\rho=1,2$ kg/m^3, $C=0,4$, $A=0,2$ m^2. Tính lực cản.
\shortans{3,89}
\loigiai{
Thay số: $F_c=\frac12\cdot1,2\cdot0,4\cdot0,2\cdot9^2=3,89$ N. Đáp án trả lời ngắn không ghi đơn vị.
}
\end{ex}

% ===== Câu 49 =====
% ID: AIQ_L10C3_FULL_1021
% Mức: TH
\begin{bt}
% ID: AIQ_L10C3_FULL_1021
% Mức: TH
Vật chuyển động trong không khí với tốc độ $10\,\text{m}$ /s. Biết lực cản $F_c=\frac{1}{2}\rho C A v^2$, $\rho=1,2$ kg/m^3, $C=0,5$, $A=0,25$ m^2. Tính lực cản.
\loigiai{
Phân tích dữ kiện, chọn định luật phù hợp. Thay số: $F_c=\frac12\cdot1,2\cdot0,5\cdot0,25\cdot10^2=7,5$ N.
}
\end{bt}

% ===== Câu 50 =====
% ID: AIQ_L10C3_FULL_1022
% Mức: TH
\begin{bt}
% ID: AIQ_L10C3_FULL_1022
% Mức: TH
Vật chuyển động trong không khí với tốc độ $11\,\text{m}$ /s. Biết lực cản $F_c=\frac{1}{2}\rho C A v^2$, $\rho=1,2$ kg/m^3, $C=0,6$, $A=0,3$ m^2. Tính lực cản.
\loigiai{
Phân tích dữ kiện, chọn định luật phù hợp. Thay số: $F_c=\frac12\cdot1,2\cdot0,6\cdot0,3\cdot11^2=13,07$ N.
}
\end{bt}

% ===== Câu 51 =====
% ID: AIQ_L10C3_FULL_1023
% Mức: VD
\begin{bt}
% ID: AIQ_L10C3_FULL_1023
% Mức: VD
Vật chuyển động trong không khí với tốc độ $12\,\text{m}$ /s. Biết lực cản $F_c=\frac{1}{2}\rho C A v^2$, $\rho=1,2$ kg/m^3, $C=0,4$, $A=0,35$ m^2. Tính lực cản.
\loigiai{
Phân tích dữ kiện, chọn định luật phù hợp. Thay số: $F_c=\frac12\cdot1,2\cdot0,4\cdot0,35\cdot12^2=12,1$ N.
}
\end{bt}

% ===== Câu 52 =====
% ID: AIQ_L10C3_FULL_1024
% Mức: VD
\begin{bt}
% ID: AIQ_L10C3_FULL_1024
% Mức: VD
Vật chuyển động trong không khí với tốc độ $5\,\text{m}$ /s. Biết lực cản $F_c=\frac{1}{2}\rho C A v^2$, $\rho=1,2$ kg/m^3, $C=0,5$, $A=0,4$ m^2. Tính lực cản.
\loigiai{
Phân tích dữ kiện, chọn định luật phù hợp. Thay số: $F_c=\frac12\cdot1,2\cdot0,5\cdot0,4\cdot5^2=3$ N.
}
\end{bt}

% ===== Câu 53 =====
% ID: AIQ_L10C3_FULL_1025
% Mức: VD
\begin{bt}
% ID: AIQ_L10C3_FULL_1025
% Mức: VD
Vật chuyển động trong không khí với tốc độ $6\,\text{m}$ /s. Biết lực cản $F_c=\frac{1}{2}\rho C A v^2$, $\rho=1,2$ kg/m^3, $C=0,6$, $A=0,2$ m^2. Tính lực cản.
\loigiai{
Phân tích dữ kiện, chọn định luật phù hợp. Thay số: $F_c=\frac12\cdot1,2\cdot0,6\cdot0,2\cdot6^2=2,59$ N.
}
\end{bt}

% ===== Câu 54 =====
% ID: AIQ_L10C3_FULL_1026
% Mức: VDC
\begin{bt}
% ID: AIQ_L10C3_FULL_1026
% Mức: VDC
Vật chuyển động trong không khí với tốc độ $7\,\text{m}$ /s. Biết lực cản $F_c=\frac{1}{2}\rho C A v^2$, $\rho=1,2$ kg/m^3, $C=0,4$, $A=0,25$ m^2. Tính lực cản.
\loigiai{
Phân tích dữ kiện, chọn định luật phù hợp. Thay số: $F_c=\frac12\cdot1,2\cdot0,4\cdot0,25\cdot7^2=2,94$ N.
}
\end{bt}

% ===== Câu 55 =====
% ID: AIQ_L10C3_FULL_1027
% Mức: NB
\dangbt{Phân tích chuyển động rơi có lực cản và vận tốc giới hạn}
\begin{ex}
% ID: AIQ_L10C3_FULL_1027
% Mức: NB
Vật khối lượng $0,35\,\text{kg}$ rơi thẳng đứng và đã đạt vận tốc giới hạn. Lấy g= $10\,\text{m}$ /s^2. Độ lớn lực cản khi đó bằng bao nhiêu?
\choice
{5,5 $N$}
{\True 3,5 $N$}
{4,2 $N$}
{2,8 $N$}
\loigiai{
Ở vận tốc giới hạn, gia tốc bằng 0 nên $F_c=P=mg=3,5$ N. Chọn B.
}
\end{ex}

% ===== Câu 56 =====
% ID: AIQ_L10C3_FULL_1028
% Mức: NB
\begin{ex}
% ID: AIQ_L10C3_FULL_1028
% Mức: NB
Vật khối lượng $0,4\,\text{kg}$ rơi thẳng đứng và đã đạt vận tốc giới hạn. Lấy g= $10\,\text{m}$ /s^2. Độ lớn lực cản khi đó bằng bao nhiêu?
\choice
{\True 4 $N$}
{4,8 $N$}
{3,2 $N$}
{6 $N$}
\loigiai{
Ở vận tốc giới hạn, gia tốc bằng 0 nên $F_c=P=mg=4$ N. Chọn A.
}
\end{ex}

% ===== Câu 57 =====
% ID: AIQ_L10C3_FULL_1029
% Mức: NB
\begin{ex}
% ID: AIQ_L10C3_FULL_1029
% Mức: NB
Vật khối lượng $0,1\,\text{kg}$ rơi thẳng đứng và đã đạt vận tốc giới hạn. Lấy g= $10\,\text{m}$ /s^2. Độ lớn lực cản khi đó bằng bao nhiêu?
\choice
{1,2 $N$}
{0,8 $N$}
{3 $N$}
{\True 1 $N$}
\loigiai{
Ở vận tốc giới hạn, gia tốc bằng 0 nên $F_c=P=mg=1$ N. Chọn D.
}
\end{ex}

% ===== Câu 58 =====
% ID: AIQ_L10C3_FULL_1030
% Mức: TH
\begin{ex}
% ID: AIQ_L10C3_FULL_1030
% Mức: TH
Vật khối lượng $0,15\,\text{kg}$ rơi thẳng đứng và đã đạt vận tốc giới hạn. Lấy g= $10\,\text{m}$ /s^2. Độ lớn lực cản khi đó bằng bao nhiêu?
\choice
{1,2 $N$}
{3,5 $N$}
{\True 1,5 $N$}
{1,8 $N$}
\loigiai{
Ở vận tốc giới hạn, gia tốc bằng 0 nên $F_c=P=mg=1,5$ N. Chọn C.
}
\end{ex}

% ===== Câu 59 =====
% ID: AIQ_L10C3_FULL_1031
% Mức: TH
\begin{ex}
% ID: AIQ_L10C3_FULL_1031
% Mức: TH
Vật khối lượng $0,2\,\text{kg}$ rơi thẳng đứng và đã đạt vận tốc giới hạn. Lấy g= $10\,\text{m}$ /s^2. Độ lớn lực cản khi đó bằng bao nhiêu?
\choice
{4 $N$}
{\True 2 $N$}
{2,4 $N$}
{1,6 $N$}
\loigiai{
Ở vận tốc giới hạn, gia tốc bằng 0 nên $F_c=P=mg=2$ N. Chọn B.
}
\end{ex}

% ===== Câu 60 =====
% ID: AIQ_L10C3_FULL_1032
% Mức: TH
\begin{ex}
% ID: AIQ_L10C3_FULL_1032
% Mức: TH
Vật khối lượng $0,25\,\text{kg}$ rơi thẳng đứng và đã đạt vận tốc giới hạn. Lấy g= $10\,\text{m}$ /s^2. Độ lớn lực cản khi đó bằng bao nhiêu?
\choice
{\True 2,5 $N$}
{3 $N$}
{2 $N$}
{4,5 $N$}
\loigiai{
Ở vận tốc giới hạn, gia tốc bằng 0 nên $F_c=P=mg=2,5$ N. Chọn A.
}
\end{ex}

% ===== Câu 61 =====
% ID: AIQ_L10C3_FULL_1033
% Mức: TH
\begin{ex}
% ID: AIQ_L10C3_FULL_1033
% Mức: TH
Vật khối lượng $0,3\,\text{kg}$ rơi thẳng đứng và đã đạt vận tốc giới hạn. Lấy g= $10\,\text{m}$ /s^2. Độ lớn lực cản khi đó bằng bao nhiêu?
\choice
{3,6 $N$}
{2,4 $N$}
{5 $N$}
{\True 3 $N$}
\loigiai{
Ở vận tốc giới hạn, gia tốc bằng 0 nên $F_c=P=mg=3$ N. Chọn D.
}
\end{ex}

% ===== Câu 62 =====
% ID: AIQ_L10C3_FULL_1034
% Mức: VD
\begin{ex}
% ID: AIQ_L10C3_FULL_1034
% Mức: VD
Vật khối lượng $0,35\,\text{kg}$ rơi thẳng đứng và đã đạt vận tốc giới hạn. Lấy g= $10\,\text{m}$ /s^2. Độ lớn lực cản khi đó bằng bao nhiêu?
\choice
{2,8 $N$}
{5,5 $N$}
{\True 3,5 $N$}
{4,2 $N$}
\loigiai{
Ở vận tốc giới hạn, gia tốc bằng 0 nên $F_c=P=mg=3,5$ N. Chọn C.
}
\end{ex}

% ===== Câu 63 =====
% ID: AIQ_L10C3_FULL_1035
% Mức: VD
\begin{ex}
% ID: AIQ_L10C3_FULL_1035
% Mức: VD
Vật khối lượng $0,4\,\text{kg}$ rơi thẳng đứng và đã đạt vận tốc giới hạn. Lấy g= $10\,\text{m}$ /s^2. Độ lớn lực cản khi đó bằng bao nhiêu?
\choice
{6 $N$}
{\True 4 $N$}
{4,8 $N$}
{3,2 $N$}
\loigiai{
Ở vận tốc giới hạn, gia tốc bằng 0 nên $F_c=P=mg=4$ N. Chọn B.
}
\end{ex}

% ===== Câu 64 =====
% ID: AIQ_L10C3_FULL_1036
% Mức: VDC
\begin{ex}
% ID: AIQ_L10C3_FULL_1036
% Mức: VDC
Vật khối lượng $0,1\,\text{kg}$ rơi thẳng đứng và đã đạt vận tốc giới hạn. Lấy g= $10\,\text{m}$ /s^2. Độ lớn lực cản khi đó bằng bao nhiêu?
\choice
{\True 1 $N$}
{1,2 $N$}
{0,8 $N$}
{3 $N$}
\loigiai{
Ở vận tốc giới hạn, gia tốc bằng 0 nên $F_c=P=mg=1$ N. Chọn A.
}
\end{ex}

% ===== Câu 65 =====
% ID: AIQ_L10C3_FULL_1037
% Mức: NB
\begin{ex}
% ID: AIQ_L10C3_FULL_1037
% Mức: NB
Vật khối lượng $0,15\,\text{kg}$ rơi thẳng đứng và đã đạt vận tốc giới hạn. Lấy g= $10\,\text{m}$ /s^2. Độ lớn lực cản khi đó bằng bao nhiêu?
\choiceTF
{\True Kết quả đúng là $1,5\,\text{N}$.}
{Kết quả bằng $3\,\text{N}$.}
{\True Khi hợp lực bằng 0, vật có gia tốc bằng 0.}
{Có thể bỏ qua đơn vị của đại lượng trong kết quả vật lí.}
\loigiai{
\begin{itemchoice}
\item (Đúng) Kết quả đúng là $1,5\,\text{N}$. (Vì): Ở vận tốc giới hạn, gia tốc bằng 0 nên $F_c=P=mg=1,5$ N. b) Sai vì giá trị hoặc chiều nêu ra không phù hợp. c) Đúng theo định luật II Newton. d) Sai vì kết quả vật lí phải kèm đơn vị thích hợp
\item (Sai) Kết quả bằng $3\,\text{N}$.
\item (Đúng) Khi hợp lực bằng 0, vật có gia tốc bằng 0.
\item (Sai) Có thể bỏ qua đơn vị của đại lượng trong kết quả vật lí.
\end{itemchoice}
}
\end{ex}

% ===== Câu 66 =====
% ID: AIQ_L10C3_FULL_1038
% Mức: NB
\begin{ex}
% ID: AIQ_L10C3_FULL_1038
% Mức: NB
Vật khối lượng $0,2\,\text{kg}$ rơi thẳng đứng và đã đạt vận tốc giới hạn. Lấy g= $10\,\text{m}$ /s^2. Độ lớn lực cản khi đó bằng bao nhiêu?
\choiceTF
{\True Kết quả đúng là $2\,\text{N}$.}
{Kết quả bằng $4\,\text{N}$.}
{\True Khi hợp lực bằng 0, vật có gia tốc bằng 0.}
{Có thể bỏ qua đơn vị của đại lượng trong kết quả vật lí.}
\loigiai{
\begin{itemchoice}
\item (Đúng) Kết quả đúng là $2\,\text{N}$. (Vì): Ở vận tốc giới hạn, gia tốc bằng 0 nên $F_c=P=mg=2$ N. b) Sai vì giá trị hoặc chiều nêu ra không phù hợp. c) Đúng theo định luật II Newton. d) Sai vì kết quả vật lí phải kèm đơn vị thích hợp
\item (Sai) Kết quả bằng $4\,\text{N}$.
\item (Đúng) Khi hợp lực bằng 0, vật có gia tốc bằng 0.
\item (Sai) Có thể bỏ qua đơn vị của đại lượng trong kết quả vật lí.
\end{itemchoice}
}
\end{ex}

% ===== Câu 67 =====
% ID: AIQ_L10C3_FULL_1039
% Mức: TH
\begin{ex}
% ID: AIQ_L10C3_FULL_1039
% Mức: TH
Vật khối lượng $0,25\,\text{kg}$ rơi thẳng đứng và đã đạt vận tốc giới hạn. Lấy g= $10\,\text{m}$ /s^2. Độ lớn lực cản khi đó bằng bao nhiêu?
\choiceTF
{\True Kết quả đúng là $2,5\,\text{N}$.}
{Kết quả bằng $5\,\text{N}$.}
{\True Khi hợp lực bằng 0, vật có gia tốc bằng 0.}
{Có thể bỏ qua đơn vị của đại lượng trong kết quả vật lí.}
\loigiai{
\begin{itemchoice}
\item (Đúng) Kết quả đúng là $2,5\,\text{N}$. (Vì): Ở vận tốc giới hạn, gia tốc bằng 0 nên $F_c=P=mg=2,5$ N. b) Sai vì giá trị hoặc chiều nêu ra không phù hợp. c) Đúng theo định luật II Newton. d) Sai vì kết quả vật lí phải kèm đơn vị thích hợp
\item (Sai) Kết quả bằng $5\,\text{N}$.
\item (Đúng) Khi hợp lực bằng 0, vật có gia tốc bằng 0.
\item (Sai) Có thể bỏ qua đơn vị của đại lượng trong kết quả vật lí.
\end{itemchoice}
}
\end{ex}

% ===== Câu 68 =====
% ID: AIQ_L10C3_FULL_1040
% Mức: TH
\begin{ex}
% ID: AIQ_L10C3_FULL_1040
% Mức: TH
Vật khối lượng $0,3\,\text{kg}$ rơi thẳng đứng và đã đạt vận tốc giới hạn. Lấy g= $10\,\text{m}$ /s^2. Độ lớn lực cản khi đó bằng bao nhiêu?
\choiceTF
{\True Kết quả đúng là $3\,\text{N}$.}
{Kết quả bằng $6\,\text{N}$.}
{\True Khi hợp lực bằng 0, vật có gia tốc bằng 0.}
{Có thể bỏ qua đơn vị của đại lượng trong kết quả vật lí.}
\loigiai{
\begin{itemchoice}
\item (Đúng) Kết quả đúng là $3\,\text{N}$. (Vì): Ở vận tốc giới hạn, gia tốc bằng 0 nên $F_c=P=mg=3$ N. b) Sai vì giá trị hoặc chiều nêu ra không phù hợp. c) Đúng theo định luật II Newton. d) Sai vì kết quả vật lí phải kèm đơn vị thích hợp
\item (Sai) Kết quả bằng $6\,\text{N}$.
\item (Đúng) Khi hợp lực bằng 0, vật có gia tốc bằng 0.
\item (Sai) Có thể bỏ qua đơn vị của đại lượng trong kết quả vật lí.
\end{itemchoice}
}
\end{ex}

% ===== Câu 69 =====
% ID: AIQ_L10C3_FULL_1041
% Mức: VD
\begin{ex}
% ID: AIQ_L10C3_FULL_1041
% Mức: VD
Vật khối lượng $0,35\,\text{kg}$ rơi thẳng đứng và đã đạt vận tốc giới hạn. Lấy g= $10\,\text{m}$ /s^2. Độ lớn lực cản khi đó bằng bao nhiêu?
\choiceTF
{\True Kết quả đúng là $3,5\,\text{N}$.}
{Kết quả bằng $7\,\text{N}$.}
{\True Khi hợp lực bằng 0, vật có gia tốc bằng 0.}
{Có thể bỏ qua đơn vị của đại lượng trong kết quả vật lí.}
\loigiai{
\begin{itemchoice}
\item (Đúng) Kết quả đúng là $3,5\,\text{N}$. (Vì): Ở vận tốc giới hạn, gia tốc bằng 0 nên $F_c=P=mg=3,5$ N. b) Sai vì giá trị hoặc chiều nêu ra không phù hợp. c) Đúng theo định luật II Newton. d) Sai vì kết quả vật lí phải kèm đơn vị thích hợp
\item (Sai) Kết quả bằng $7\,\text{N}$.
\item (Đúng) Khi hợp lực bằng 0, vật có gia tốc bằng 0.
\item (Sai) Có thể bỏ qua đơn vị của đại lượng trong kết quả vật lí.
\end{itemchoice}
}
\end{ex}

% ===== Câu 70 =====
% ID: AIQ_L10C3_FULL_1042
% Mức: TH
\begin{ex}
% ID: AIQ_L10C3_FULL_1042
% Mức: TH
Vật khối lượng $0,4\,\text{kg}$ rơi thẳng đứng và đã đạt vận tốc giới hạn. Lấy g= $10\,\text{m}$ /s^2. Độ lớn lực cản khi đó bằng bao nhiêu?
\shortans{4}
\loigiai{
Ở vận tốc giới hạn, gia tốc bằng 0 nên $F_c=P=mg=4$ N. Đáp án trả lời ngắn không ghi đơn vị.
}
\end{ex}

% ===== Câu 71 =====
% ID: AIQ_L10C3_FULL_1043
% Mức: TH
\begin{ex}
% ID: AIQ_L10C3_FULL_1043
% Mức: TH
Vật khối lượng $0,1\,\text{kg}$ rơi thẳng đứng và đã đạt vận tốc giới hạn. Lấy g= $10\,\text{m}$ /s^2. Độ lớn lực cản khi đó bằng bao nhiêu?
\shortans{1}
\loigiai{
Ở vận tốc giới hạn, gia tốc bằng 0 nên $F_c=P=mg=1$ N. Đáp án trả lời ngắn không ghi đơn vị.
}
\end{ex}

% ===== Câu 72 =====
% ID: AIQ_L10C3_FULL_1044
% Mức: VD
\begin{ex}
% ID: AIQ_L10C3_FULL_1044
% Mức: VD
Vật khối lượng $0,15\,\text{kg}$ rơi thẳng đứng và đã đạt vận tốc giới hạn. Lấy g= $10\,\text{m}$ /s^2. Độ lớn lực cản khi đó bằng bao nhiêu?
\shortans{1,5}
\loigiai{
Ở vận tốc giới hạn, gia tốc bằng 0 nên $F_c=P=mg=1,5$ N. Đáp án trả lời ngắn không ghi đơn vị.
}
\end{ex}

% ===== Câu 73 =====
% ID: AIQ_L10C3_FULL_1045
% Mức: VD
\begin{ex}
% ID: AIQ_L10C3_FULL_1045
% Mức: VD
Vật khối lượng $0,2\,\text{kg}$ rơi thẳng đứng và đã đạt vận tốc giới hạn. Lấy g= $10\,\text{m}$ /s^2. Độ lớn lực cản khi đó bằng bao nhiêu?
\shortans{2}
\loigiai{
Ở vận tốc giới hạn, gia tốc bằng 0 nên $F_c=P=mg=2$ N. Đáp án trả lời ngắn không ghi đơn vị.
}
\end{ex}

% ===== Câu 74 =====
% ID: AIQ_L10C3_FULL_1046
% Mức: VD
\begin{ex}
% ID: AIQ_L10C3_FULL_1046
% Mức: VD
Vật khối lượng $0,25\,\text{kg}$ rơi thẳng đứng và đã đạt vận tốc giới hạn. Lấy g= $10\,\text{m}$ /s^2. Độ lớn lực cản khi đó bằng bao nhiêu?
\shortans{2,5}
\loigiai{
Ở vận tốc giới hạn, gia tốc bằng 0 nên $F_c=P=mg=2,5$ N. Đáp án trả lời ngắn không ghi đơn vị.
}
\end{ex}

% ===== Câu 75 =====
% ID: AIQ_L10C3_FULL_1047
% Mức: VDC
\begin{ex}
% ID: AIQ_L10C3_FULL_1047
% Mức: VDC
Vật khối lượng $0,3\,\text{kg}$ rơi thẳng đứng và đã đạt vận tốc giới hạn. Lấy g= $10\,\text{m}$ /s^2. Độ lớn lực cản khi đó bằng bao nhiêu?
\shortans{3}
\loigiai{
Ở vận tốc giới hạn, gia tốc bằng 0 nên $F_c=P=mg=3$ N. Đáp án trả lời ngắn không ghi đơn vị.
}
\end{ex}

% ===== Câu 76 =====
% ID: AIQ_L10C3_FULL_1048
% Mức: TH
\begin{bt}
% ID: AIQ_L10C3_FULL_1048
% Mức: TH
Vật khối lượng $0,35\,\text{kg}$ rơi thẳng đứng và đã đạt vận tốc giới hạn. Lấy g= $10\,\text{m}$ /s^2. Độ lớn lực cản khi đó bằng bao nhiêu?
\loigiai{
Phân tích dữ kiện, chọn định luật phù hợp. Ở vận tốc giới hạn, gia tốc bằng 0 nên $F_c=P=mg=3,5$ N.
}
\end{bt}

% ===== Câu 77 =====
% ID: AIQ_L10C3_FULL_1049
% Mức: TH
\begin{bt}
% ID: AIQ_L10C3_FULL_1049
% Mức: TH
Vật khối lượng $0,4\,\text{kg}$ rơi thẳng đứng và đã đạt vận tốc giới hạn. Lấy g= $10\,\text{m}$ /s^2. Độ lớn lực cản khi đó bằng bao nhiêu?
\loigiai{
Phân tích dữ kiện, chọn định luật phù hợp. Ở vận tốc giới hạn, gia tốc bằng 0 nên $F_c=P=mg=4$ N.
}
\end{bt}

% ===== Câu 78 =====
% ID: AIQ_L10C3_FULL_1050
% Mức: VD
\begin{bt}
% ID: AIQ_L10C3_FULL_1050
% Mức: VD
Vật khối lượng $0,1\,\text{kg}$ rơi thẳng đứng và đã đạt vận tốc giới hạn. Lấy g= $10\,\text{m}$ /s^2. Độ lớn lực cản khi đó bằng bao nhiêu?
\loigiai{
Phân tích dữ kiện, chọn định luật phù hợp. Ở vận tốc giới hạn, gia tốc bằng 0 nên $F_c=P=mg=1$ N.
}
\end{bt}

% ===== Câu 79 =====
% ID: AIQ_L10C3_FULL_1051
% Mức: VD
\begin{bt}
% ID: AIQ_L10C3_FULL_1051
% Mức: VD
Vật khối lượng $0,15\,\text{kg}$ rơi thẳng đứng và đã đạt vận tốc giới hạn. Lấy g= $10\,\text{m}$ /s^2. Độ lớn lực cản khi đó bằng bao nhiêu?
\loigiai{
Phân tích dữ kiện, chọn định luật phù hợp. Ở vận tốc giới hạn, gia tốc bằng 0 nên $F_c=P=mg=1,5$ N.
}
\end{bt}

% ===== Câu 80 =====
% ID: AIQ_L10C3_FULL_1052
% Mức: VD
\begin{bt}
% ID: AIQ_L10C3_FULL_1052
% Mức: VD
Vật khối lượng $0,2\,\text{kg}$ rơi thẳng đứng và đã đạt vận tốc giới hạn. Lấy g= $10\,\text{m}$ /s^2. Độ lớn lực cản khi đó bằng bao nhiêu?
\loigiai{
Phân tích dữ kiện, chọn định luật phù hợp. Ở vận tốc giới hạn, gia tốc bằng 0 nên $F_c=P=mg=2$ N.
}
\end{bt}

% ===== Câu 81 =====
% ID: AIQ_L10C3_FULL_1053
% Mức: VDC
\begin{bt}
% ID: AIQ_L10C3_FULL_1053
% Mức: VDC
Vật khối lượng $0,25\,\text{kg}$ rơi thẳng đứng và đã đạt vận tốc giới hạn. Lấy g= $10\,\text{m}$ /s^2. Độ lớn lực cản khi đó bằng bao nhiêu?
\loigiai{
Phân tích dữ kiện, chọn định luật phù hợp. Ở vận tốc giới hạn, gia tốc bằng 0 nên $F_c=P=mg=2,5$ N.
}
\end{bt}

% ===== Câu 82 =====
% ID: AIQ_L10C3_FULL_1054
% Mức: NB
\dangbt{Nhận biết và giải thích lực nâng của chất lưu}
\begin{ex}
% ID: AIQ_L10C3_FULL_1054
% Mức: NB
Cánh máy bay tạo lực nâng chủ yếu nhờ yếu tố nào?
\choice
{\True Chênh lệch áp suất giữa hai mặt cánh}
{Khối lượng cánh tăng}
{Trọng lực biến mất}
{Lực cản luôn hướng lên}
\loigiai{
Dòng khí quanh cánh tạo chênh lệch áp suất; hợp lực áp suất có thành phần hướng lên. Chọn A.
}
\end{ex}

% ===== Câu 83 =====
% ID: AIQ_L10C3_FULL_1055
% Mức: NB
\begin{ex}
% ID: AIQ_L10C3_FULL_1055
% Mức: NB
Cánh máy bay tạo lực nâng chủ yếu nhờ yếu tố nào?
\choice
{\True Chênh lệch áp suất giữa hai mặt cánh}
{Khối lượng cánh tăng}
{Trọng lực biến mất}
{Lực cản luôn hướng lên}
\loigiai{
Dòng khí quanh cánh tạo chênh lệch áp suất; hợp lực áp suất có thành phần hướng lên. Chọn A.
}
\end{ex}

% ===== Câu 84 =====
% ID: AIQ_L10C3_FULL_1056
% Mức: NB
\begin{ex}
% ID: AIQ_L10C3_FULL_1056
% Mức: NB
Cánh máy bay tạo lực nâng chủ yếu nhờ yếu tố nào?
\choice
{\True Chênh lệch áp suất giữa hai mặt cánh}
{Khối lượng cánh tăng}
{Trọng lực biến mất}
{Lực cản luôn hướng lên}
\loigiai{
Dòng khí quanh cánh tạo chênh lệch áp suất; hợp lực áp suất có thành phần hướng lên. Chọn A.
}
\end{ex}

% ===== Câu 85 =====
% ID: AIQ_L10C3_FULL_1057
% Mức: TH
\begin{ex}
% ID: AIQ_L10C3_FULL_1057
% Mức: TH
Cánh máy bay tạo lực nâng chủ yếu nhờ yếu tố nào?
\choice
{\True Chênh lệch áp suất giữa hai mặt cánh}
{Khối lượng cánh tăng}
{Trọng lực biến mất}
{Lực cản luôn hướng lên}
\loigiai{
Dòng khí quanh cánh tạo chênh lệch áp suất; hợp lực áp suất có thành phần hướng lên. Chọn A.
}
\end{ex}

% ===== Câu 86 =====
% ID: AIQ_L10C3_FULL_1058
% Mức: TH
\begin{ex}
% ID: AIQ_L10C3_FULL_1058
% Mức: TH
Cánh máy bay tạo lực nâng chủ yếu nhờ yếu tố nào?
\choice
{\True Chênh lệch áp suất giữa hai mặt cánh}
{Khối lượng cánh tăng}
{Trọng lực biến mất}
{Lực cản luôn hướng lên}
\loigiai{
Dòng khí quanh cánh tạo chênh lệch áp suất; hợp lực áp suất có thành phần hướng lên. Chọn A.
}
\end{ex}

% ===== Câu 87 =====
% ID: AIQ_L10C3_FULL_1059
% Mức: TH
\begin{ex}
% ID: AIQ_L10C3_FULL_1059
% Mức: TH
Cánh máy bay tạo lực nâng chủ yếu nhờ yếu tố nào?
\choice
{\True Chênh lệch áp suất giữa hai mặt cánh}
{Khối lượng cánh tăng}
{Trọng lực biến mất}
{Lực cản luôn hướng lên}
\loigiai{
Dòng khí quanh cánh tạo chênh lệch áp suất; hợp lực áp suất có thành phần hướng lên. Chọn A.
}
\end{ex}

% ===== Câu 88 =====
% ID: AIQ_L10C3_FULL_1060
% Mức: TH
\begin{ex}
% ID: AIQ_L10C3_FULL_1060
% Mức: TH
Cánh máy bay tạo lực nâng chủ yếu nhờ yếu tố nào?
\choice
{\True Chênh lệch áp suất giữa hai mặt cánh}
{Khối lượng cánh tăng}
{Trọng lực biến mất}
{Lực cản luôn hướng lên}
\loigiai{
Dòng khí quanh cánh tạo chênh lệch áp suất; hợp lực áp suất có thành phần hướng lên. Chọn A.
}
\end{ex}

% ===== Câu 89 =====
% ID: AIQ_L10C3_FULL_1061
% Mức: VD
\begin{ex}
% ID: AIQ_L10C3_FULL_1061
% Mức: VD
Cánh máy bay tạo lực nâng chủ yếu nhờ yếu tố nào?
\choice
{\True Chênh lệch áp suất giữa hai mặt cánh}
{Khối lượng cánh tăng}
{Trọng lực biến mất}
{Lực cản luôn hướng lên}
\loigiai{
Dòng khí quanh cánh tạo chênh lệch áp suất; hợp lực áp suất có thành phần hướng lên. Chọn A.
}
\end{ex}

% ===== Câu 90 =====
% ID: AIQ_L10C3_FULL_1062
% Mức: VD
\begin{ex}
% ID: AIQ_L10C3_FULL_1062
% Mức: VD
Cánh máy bay tạo lực nâng chủ yếu nhờ yếu tố nào?
\choice
{\True Chênh lệch áp suất giữa hai mặt cánh}
{Khối lượng cánh tăng}
{Trọng lực biến mất}
{Lực cản luôn hướng lên}
\loigiai{
Dòng khí quanh cánh tạo chênh lệch áp suất; hợp lực áp suất có thành phần hướng lên. Chọn A.
}
\end{ex}

% ===== Câu 91 =====
% ID: AIQ_L10C3_FULL_1063
% Mức: VDC
\begin{ex}
% ID: AIQ_L10C3_FULL_1063
% Mức: VDC
Cánh máy bay tạo lực nâng chủ yếu nhờ yếu tố nào?
\choice
{\True Chênh lệch áp suất giữa hai mặt cánh}
{Khối lượng cánh tăng}
{Trọng lực biến mất}
{Lực cản luôn hướng lên}
\loigiai{
Dòng khí quanh cánh tạo chênh lệch áp suất; hợp lực áp suất có thành phần hướng lên. Chọn A.
}
\end{ex}

% ===== Câu 92 =====
% ID: AIQ_L10C3_FULL_1064
% Mức: NB
\begin{ex}
% ID: AIQ_L10C3_FULL_1064
% Mức: NB
Cánh máy bay tạo lực nâng chủ yếu nhờ yếu tố nào?
\choiceTF
{\True Lực cản phụ thuộc chuyển động tương đối giữa vật và chất lưu.}
{Lực cản luôn cùng chiều chuyển động của vật.}
{\True Khi hợp lực bằng 0, vật có gia tốc bằng 0.}
{Có thể bỏ qua đơn vị của đại lượng trong kết quả vật lí.}
\loigiai{
\begin{itemchoice}
\item (Đúng) Lực cản phụ thuộc chuyển động tương đối giữa vật và chất lưu. (Vì): òng khí quanh cánh tạo chênh lệch áp suất; hợp lực áp suất có thành phần hướng lên. b) Sai vì giá trị hoặc chiều nêu ra không phù hợp. c) Đúng theo định luật II Newton. d) Sai vì kết quả vật lí phải kèm đơn vị thích hợp
\item (Sai) Lực cản luôn cùng chiều chuyển động của vật.
\item (Đúng) Khi hợp lực bằng 0, vật có gia tốc bằng 0.
\item (Sai) Có thể bỏ qua đơn vị của đại lượng trong kết quả vật lí.
\end{itemchoice}
}
\end{ex}

% ===== Câu 93 =====
% ID: AIQ_L10C3_FULL_1065
% Mức: NB
\begin{ex}
% ID: AIQ_L10C3_FULL_1065
% Mức: NB
Cánh máy bay tạo lực nâng chủ yếu nhờ yếu tố nào?
\choiceTF
{\True Lực cản phụ thuộc chuyển động tương đối giữa vật và chất lưu.}
{Lực cản luôn cùng chiều chuyển động của vật.}
{\True Khi hợp lực bằng 0, vật có gia tốc bằng 0.}
{Có thể bỏ qua đơn vị của đại lượng trong kết quả vật lí.}
\loigiai{
\begin{itemchoice}
\item (Đúng) Lực cản phụ thuộc chuyển động tương đối giữa vật và chất lưu. (Vì): òng khí quanh cánh tạo chênh lệch áp suất; hợp lực áp suất có thành phần hướng lên. b) Sai vì giá trị hoặc chiều nêu ra không phù hợp. c) Đúng theo định luật II Newton. d) Sai vì kết quả vật lí phải kèm đơn vị thích hợp
\item (Sai) Lực cản luôn cùng chiều chuyển động của vật.
\item (Đúng) Khi hợp lực bằng 0, vật có gia tốc bằng 0.
\item (Sai) Có thể bỏ qua đơn vị của đại lượng trong kết quả vật lí.
\end{itemchoice}
}
\end{ex}

% ===== Câu 94 =====
% ID: AIQ_L10C3_FULL_1066
% Mức: TH
\begin{ex}
% ID: AIQ_L10C3_FULL_1066
% Mức: TH
Cánh máy bay tạo lực nâng chủ yếu nhờ yếu tố nào?
\choiceTF
{\True Lực cản phụ thuộc chuyển động tương đối giữa vật và chất lưu.}
{Lực cản luôn cùng chiều chuyển động của vật.}
{\True Khi hợp lực bằng 0, vật có gia tốc bằng 0.}
{Có thể bỏ qua đơn vị của đại lượng trong kết quả vật lí.}
\loigiai{
\begin{itemchoice}
\item (Đúng) Lực cản phụ thuộc chuyển động tương đối giữa vật và chất lưu. (Vì): òng khí quanh cánh tạo chênh lệch áp suất; hợp lực áp suất có thành phần hướng lên. b) Sai vì giá trị hoặc chiều nêu ra không phù hợp. c) Đúng theo định luật II Newton. d) Sai vì kết quả vật lí phải kèm đơn vị thích hợp
\item (Sai) Lực cản luôn cùng chiều chuyển động của vật.
\item (Đúng) Khi hợp lực bằng 0, vật có gia tốc bằng 0.
\item (Sai) Có thể bỏ qua đơn vị của đại lượng trong kết quả vật lí.
\end{itemchoice}
}
\end{ex}

% ===== Câu 95 =====
% ID: AIQ_L10C3_FULL_1067
% Mức: TH
\begin{ex}
% ID: AIQ_L10C3_FULL_1067
% Mức: TH
Cánh máy bay tạo lực nâng chủ yếu nhờ yếu tố nào?
\choiceTF
{\True Lực cản phụ thuộc chuyển động tương đối giữa vật và chất lưu.}
{Lực cản luôn cùng chiều chuyển động của vật.}
{\True Khi hợp lực bằng 0, vật có gia tốc bằng 0.}
{Có thể bỏ qua đơn vị của đại lượng trong kết quả vật lí.}
\loigiai{
\begin{itemchoice}
\item (Đúng) Lực cản phụ thuộc chuyển động tương đối giữa vật và chất lưu. (Vì): òng khí quanh cánh tạo chênh lệch áp suất; hợp lực áp suất có thành phần hướng lên. b) Sai vì giá trị hoặc chiều nêu ra không phù hợp. c) Đúng theo định luật II Newton. d) Sai vì kết quả vật lí phải kèm đơn vị thích hợp
\item (Sai) Lực cản luôn cùng chiều chuyển động của vật.
\item (Đúng) Khi hợp lực bằng 0, vật có gia tốc bằng 0.
\item (Sai) Có thể bỏ qua đơn vị của đại lượng trong kết quả vật lí.
\end{itemchoice}
}
\end{ex}

% ===== Câu 96 =====
% ID: AIQ_L10C3_FULL_1068
% Mức: VD
\begin{ex}
% ID: AIQ_L10C3_FULL_1068
% Mức: VD
Cánh máy bay tạo lực nâng chủ yếu nhờ yếu tố nào?
\choiceTF
{\True Lực cản phụ thuộc chuyển động tương đối giữa vật và chất lưu.}
{Lực cản luôn cùng chiều chuyển động của vật.}
{\True Khi hợp lực bằng 0, vật có gia tốc bằng 0.}
{Có thể bỏ qua đơn vị của đại lượng trong kết quả vật lí.}
\loigiai{
\begin{itemchoice}
\item (Đúng) Lực cản phụ thuộc chuyển động tương đối giữa vật và chất lưu. (Vì): òng khí quanh cánh tạo chênh lệch áp suất; hợp lực áp suất có thành phần hướng lên. b) Sai vì giá trị hoặc chiều nêu ra không phù hợp. c) Đúng theo định luật II Newton. d) Sai vì kết quả vật lí phải kèm đơn vị thích hợp
\item (Sai) Lực cản luôn cùng chiều chuyển động của vật.
\item (Đúng) Khi hợp lực bằng 0, vật có gia tốc bằng 0.
\item (Sai) Có thể bỏ qua đơn vị của đại lượng trong kết quả vật lí.
\end{itemchoice}
}
\end{ex}

% ===== Câu 97 =====
% ID: AIQ_L10C3_FULL_1069
% Mức: TH
\begin{ex}
% ID: AIQ_L10C3_FULL_1069
% Mức: TH
Cánh máy bay tạo lực nâng chủ yếu nhờ yếu tố nào?
\shortans{Chênh}
\loigiai{
Dòng khí quanh cánh tạo chênh lệch áp suất; hợp lực áp suất có thành phần hướng lên. Đáp án trả lời ngắn không ghi đơn vị.
}
\end{ex}

% ===== Câu 98 =====
% ID: AIQ_L10C3_FULL_1070
% Mức: TH
\begin{ex}
% ID: AIQ_L10C3_FULL_1070
% Mức: TH
Cánh máy bay tạo lực nâng chủ yếu nhờ yếu tố nào?
\shortans{Chênh}
\loigiai{
Dòng khí quanh cánh tạo chênh lệch áp suất; hợp lực áp suất có thành phần hướng lên. Đáp án trả lời ngắn không ghi đơn vị.
}
\end{ex}

% ===== Câu 99 =====
% ID: AIQ_L10C3_FULL_1071
% Mức: VD
\begin{ex}
% ID: AIQ_L10C3_FULL_1071
% Mức: VD
Cánh máy bay tạo lực nâng chủ yếu nhờ yếu tố nào?
\shortans{Chênh}
\loigiai{
Dòng khí quanh cánh tạo chênh lệch áp suất; hợp lực áp suất có thành phần hướng lên. Đáp án trả lời ngắn không ghi đơn vị.
}
\end{ex}

% ===== Câu 100 =====
% ID: AIQ_L10C3_FULL_1072
% Mức: VD
\begin{ex}
% ID: AIQ_L10C3_FULL_1072
% Mức: VD
Cánh máy bay tạo lực nâng chủ yếu nhờ yếu tố nào?
\shortans{Chênh}
\loigiai{
Dòng khí quanh cánh tạo chênh lệch áp suất; hợp lực áp suất có thành phần hướng lên. Đáp án trả lời ngắn không ghi đơn vị.
}
\end{ex}

% ===== Câu 101 =====
% ID: AIQ_L10C3_FULL_1073
% Mức: VD
\begin{ex}
% ID: AIQ_L10C3_FULL_1073
% Mức: VD
Cánh máy bay tạo lực nâng chủ yếu nhờ yếu tố nào?
\shortans{Chênh}
\loigiai{
Dòng khí quanh cánh tạo chênh lệch áp suất; hợp lực áp suất có thành phần hướng lên. Đáp án trả lời ngắn không ghi đơn vị.
}
\end{ex}

% ===== Câu 102 =====
% ID: AIQ_L10C3_FULL_1074
% Mức: VDC
\begin{ex}
% ID: AIQ_L10C3_FULL_1074
% Mức: VDC
Cánh máy bay tạo lực nâng chủ yếu nhờ yếu tố nào?
\shortans{Chênh}
\loigiai{
Dòng khí quanh cánh tạo chênh lệch áp suất; hợp lực áp suất có thành phần hướng lên. Đáp án trả lời ngắn không ghi đơn vị.
}
\end{ex}

% ===== Câu 103 =====
% ID: AIQ_L10C3_FULL_1075
% Mức: TH
\begin{bt}
% ID: AIQ_L10C3_FULL_1075
% Mức: TH
Cánh máy bay tạo lực nâng chủ yếu nhờ yếu tố nào?
\loigiai{
Phân tích dữ kiện, chọn định luật phù hợp. Dòng khí quanh cánh tạo chênh lệch áp suất; hợp lực áp suất có thành phần hướng lên.
}
\end{bt}

% ===== Câu 104 =====
% ID: AIQ_L10C3_FULL_1076
% Mức: TH
\begin{bt}
% ID: AIQ_L10C3_FULL_1076
% Mức: TH
Cánh máy bay tạo lực nâng chủ yếu nhờ yếu tố nào?
\loigiai{
Phân tích dữ kiện, chọn định luật phù hợp. Dòng khí quanh cánh tạo chênh lệch áp suất; hợp lực áp suất có thành phần hướng lên.
}
\end{bt}

% ===== Câu 105 =====
% ID: AIQ_L10C3_FULL_1077
% Mức: VD
\begin{bt}
% ID: AIQ_L10C3_FULL_1077
% Mức: VD
Cánh máy bay tạo lực nâng chủ yếu nhờ yếu tố nào?
\loigiai{
Phân tích dữ kiện, chọn định luật phù hợp. Dòng khí quanh cánh tạo chênh lệch áp suất; hợp lực áp suất có thành phần hướng lên.
}
\end{bt}

% ===== Câu 106 =====
% ID: AIQ_L10C3_FULL_1078
% Mức: VD
\begin{bt}
% ID: AIQ_L10C3_FULL_1078
% Mức: VD
Cánh máy bay tạo lực nâng chủ yếu nhờ yếu tố nào?
\loigiai{
Phân tích dữ kiện, chọn định luật phù hợp. Dòng khí quanh cánh tạo chênh lệch áp suất; hợp lực áp suất có thành phần hướng lên.
}
\end{bt}

% ===== Câu 107 =====
% ID: AIQ_L10C3_FULL_1079
% Mức: VD
\begin{bt}
% ID: AIQ_L10C3_FULL_1079
% Mức: VD
Cánh máy bay tạo lực nâng chủ yếu nhờ yếu tố nào?
\loigiai{
Phân tích dữ kiện, chọn định luật phù hợp. Dòng khí quanh cánh tạo chênh lệch áp suất; hợp lực áp suất có thành phần hướng lên.
}
\end{bt}

% ===== Câu 108 =====
% ID: AIQ_L10C3_FULL_1080
% Mức: VDC
\begin{bt}
% ID: AIQ_L10C3_FULL_1080
% Mức: VDC
Cánh máy bay tạo lực nâng chủ yếu nhờ yếu tố nào?
\loigiai{
Phân tích dữ kiện, chọn định luật phù hợp. Dòng khí quanh cánh tạo chênh lệch áp suất; hợp lực áp suất có thành phần hướng lên.
}
\end{bt}

% ===== Câu 109 =====
% ID: AIQ_L10C3_FULL_1081
% Mức: NB
\dangbt{Tính lực nâng do chênh lệch áp suất}
\begin{ex}
% ID: AIQ_L10C3_FULL_1081
% Mức: NB
Hai mặt cánh có độ chênh áp $250\,\text{Pa}$ trên diện tích $0,5\,\text{m}$ ^2. Coi lực nâng $F=\Delta pA$. Tính lực nâng.
\choice
{150 $N$}
{100 $N$}
{127 $N$}
{\True 125 $N$}
\loigiai{
$F=\Delta pA=250\cdot0,5=125$ N. Chọn D.
}
\end{ex}

% ===== Câu 110 =====
% ID: AIQ_L10C3_FULL_1082
% Mức: NB
\begin{ex}
% ID: AIQ_L10C3_FULL_1082
% Mức: NB
Hai mặt cánh có độ chênh áp $300\,\text{Pa}$ trên diện tích $0,6\,\text{m}$ ^2. Coi lực nâng $F=\Delta pA$. Tính lực nâng.
\choice
{144 $N$}
{182 $N$}
{\True 180 $N$}
{216 $N$}
\loigiai{
$F=\Delta pA=300\cdot0,6=180$ N. Chọn C.
}
\end{ex}

% ===== Câu 111 =====
% ID: AIQ_L10C3_FULL_1083
% Mức: NB
\begin{ex}
% ID: AIQ_L10C3_FULL_1083
% Mức: NB
Hai mặt cánh có độ chênh áp $350\,\text{Pa}$ trên diện tích $0,7\,\text{m}$ ^2. Coi lực nâng $F=\Delta pA$. Tính lực nâng.
\choice
{247 $N$}
{\True 245 $N$}
{294 $N$}
{196 $N$}
\loigiai{
$F=\Delta pA=350\cdot0,7=245$ N. Chọn B.
}
\end{ex}

% ===== Câu 112 =====
% ID: AIQ_L10C3_FULL_1084
% Mức: TH
\begin{ex}
% ID: AIQ_L10C3_FULL_1084
% Mức: TH
Hai mặt cánh có độ chênh áp $400\,\text{Pa}$ trên diện tích $0,8\,\text{m}$ ^2. Coi lực nâng $F=\Delta pA$. Tính lực nâng.
\choice
{\True 320 $N$}
{384 $N$}
{256 $N$}
{322 $N$}
\loigiai{
$F=\Delta pA=400\cdot0,8=320$ N. Chọn A.
}
\end{ex}

% ===== Câu 113 =====
% ID: AIQ_L10C3_FULL_1085
% Mức: TH
\begin{ex}
% ID: AIQ_L10C3_FULL_1085
% Mức: TH
Hai mặt cánh có độ chênh áp $450\,\text{Pa}$ trên diện tích $0,9\,\text{m}$ ^2. Coi lực nâng $F=\Delta pA$. Tính lực nâng.
\choice
{486 $N$}
{324 $N$}
{407 $N$}
{\True 405 $N$}
\loigiai{
$F=\Delta pA=450\cdot0,9=405$ N. Chọn D.
}
\end{ex}

% ===== Câu 114 =====
% ID: AIQ_L10C3_FULL_1086
% Mức: TH
\begin{ex}
% ID: AIQ_L10C3_FULL_1086
% Mức: TH
Hai mặt cánh có độ chênh áp $500\,\text{Pa}$ trên diện tích $0,4\,\text{m}$ ^2. Coi lực nâng $F=\Delta pA$. Tính lực nâng.
\choice
{160 $N$}
{202 $N$}
{\True 200 $N$}
{240 $N$}
\loigiai{
$F=\Delta pA=500\cdot0,4=200$ N. Chọn C.
}
\end{ex}

% ===== Câu 115 =====
% ID: AIQ_L10C3_FULL_1087
% Mức: TH
\begin{ex}
% ID: AIQ_L10C3_FULL_1087
% Mức: TH
Hai mặt cánh có độ chênh áp $550\,\text{Pa}$ trên diện tích $0,5\,\text{m}$ ^2. Coi lực nâng $F=\Delta pA$. Tính lực nâng.
\choice
{277 $N$}
{\True 275 $N$}
{330 $N$}
{220 $N$}
\loigiai{
$F=\Delta pA=550\cdot0,5=275$ N. Chọn B.
}
\end{ex}

% ===== Câu 116 =====
% ID: AIQ_L10C3_FULL_1088
% Mức: VD
\begin{ex}
% ID: AIQ_L10C3_FULL_1088
% Mức: VD
Hai mặt cánh có độ chênh áp $200\,\text{Pa}$ trên diện tích $0,6\,\text{m}$ ^2. Coi lực nâng $F=\Delta pA$. Tính lực nâng.
\choice
{\True 120 $N$}
{144 $N$}
{96 $N$}
{122 $N$}
\loigiai{
$F=\Delta pA=200\cdot0,6=120$ N. Chọn A.
}
\end{ex}

% ===== Câu 117 =====
% ID: AIQ_L10C3_FULL_1089
% Mức: VD
\begin{ex}
% ID: AIQ_L10C3_FULL_1089
% Mức: VD
Hai mặt cánh có độ chênh áp $250\,\text{Pa}$ trên diện tích $0,7\,\text{m}$ ^2. Coi lực nâng $F=\Delta pA$. Tính lực nâng.
\choice
{210 $N$}
{140 $N$}
{177 $N$}
{\True 175 $N$}
\loigiai{
$F=\Delta pA=250\cdot0,7=175$ N. Chọn D.
}
\end{ex}

% ===== Câu 118 =====
% ID: AIQ_L10C3_FULL_1090
% Mức: VDC
\begin{ex}
% ID: AIQ_L10C3_FULL_1090
% Mức: VDC
Hai mặt cánh có độ chênh áp $300\,\text{Pa}$ trên diện tích $0,8\,\text{m}$ ^2. Coi lực nâng $F=\Delta pA$. Tính lực nâng.
\choice
{192 $N$}
{242 $N$}
{\True 240 $N$}
{288 $N$}
\loigiai{
$F=\Delta pA=300\cdot0,8=240$ N. Chọn C.
}
\end{ex}

% ===== Câu 119 =====
% ID: AIQ_L10C3_FULL_1091
% Mức: NB
\begin{ex}
% ID: AIQ_L10C3_FULL_1091
% Mức: NB
Hai mặt cánh có độ chênh áp $350\,\text{Pa}$ trên diện tích $0,9\,\text{m}$ ^2. Coi lực nâng $F=\Delta pA$. Tính lực nâng.
\choiceTF
{\True Kết quả đúng là $315\,\text{N}$.}
{Kết quả bằng $630\,\text{N}$.}
{\True Khi hợp lực bằng 0, vật có gia tốc bằng 0.}
{Có thể bỏ qua đơn vị của đại lượng trong kết quả vật lí.}
\loigiai{
\begin{itemchoice}
\item (Đúng) Kết quả đúng là $315\,\text{N}$. (Vì): $F=\Delta pA=350\cdot0,9=315$ N. b) Sai vì giá trị hoặc chiều nêu ra không phù hợp. c) Đúng theo định luật II Newton. d) Sai vì kết quả vật lí phải kèm đơn vị thích hợp
\item (Sai) Kết quả bằng $630\,\text{N}$.
\item (Đúng) Khi hợp lực bằng 0, vật có gia tốc bằng 0.
\item (Sai) Có thể bỏ qua đơn vị của đại lượng trong kết quả vật lí.
\end{itemchoice}
}
\end{ex}

% ===== Câu 120 =====
% ID: AIQ_L10C3_FULL_1092
% Mức: NB
\begin{ex}
% ID: AIQ_L10C3_FULL_1092
% Mức: NB
Hai mặt cánh có độ chênh áp $400\,\text{Pa}$ trên diện tích $0,4\,\text{m}$ ^2. Coi lực nâng $F=\Delta pA$. Tính lực nâng.
\choiceTF
{\True Kết quả đúng là $160\,\text{N}$.}
{Kết quả bằng $320\,\text{N}$.}
{\True Khi hợp lực bằng 0, vật có gia tốc bằng 0.}
{Có thể bỏ qua đơn vị của đại lượng trong kết quả vật lí.}
\loigiai{
\begin{itemchoice}
\item (Đúng) Kết quả đúng là $160\,\text{N}$. (Vì): $F=\Delta pA=400\cdot0,4=160$ N. b) Sai vì giá trị hoặc chiều nêu ra không phù hợp. c) Đúng theo định luật II Newton. d) Sai vì kết quả vật lí phải kèm đơn vị thích hợp
\item (Sai) Kết quả bằng $320\,\text{N}$.
\item (Đúng) Khi hợp lực bằng 0, vật có gia tốc bằng 0.
\item (Sai) Có thể bỏ qua đơn vị của đại lượng trong kết quả vật lí.
\end{itemchoice}
}
\end{ex}

% ===== Câu 121 =====
% ID: AIQ_L10C3_FULL_1093
% Mức: TH
\begin{ex}
% ID: AIQ_L10C3_FULL_1093
% Mức: TH
Hai mặt cánh có độ chênh áp $450\,\text{Pa}$ trên diện tích $0,5\,\text{m}$ ^2. Coi lực nâng $F=\Delta pA$. Tính lực nâng.
\choiceTF
{\True Kết quả đúng là $225\,\text{N}$.}
{Kết quả bằng $450\,\text{N}$.}
{\True Khi hợp lực bằng 0, vật có gia tốc bằng 0.}
{Có thể bỏ qua đơn vị của đại lượng trong kết quả vật lí.}
\loigiai{
\begin{itemchoice}
\item (Đúng) Kết quả đúng là $225\,\text{N}$. (Vì): $F=\Delta pA=450\cdot0,5=225$ N. b) Sai vì giá trị hoặc chiều nêu ra không phù hợp. c) Đúng theo định luật II Newton. d) Sai vì kết quả vật lí phải kèm đơn vị thích hợp
\item (Sai) Kết quả bằng $450\,\text{N}$.
\item (Đúng) Khi hợp lực bằng 0, vật có gia tốc bằng 0.
\item (Sai) Có thể bỏ qua đơn vị của đại lượng trong kết quả vật lí.
\end{itemchoice}
}
\end{ex}

% ===== Câu 122 =====
% ID: AIQ_L10C3_FULL_1094
% Mức: TH
\begin{ex}
% ID: AIQ_L10C3_FULL_1094
% Mức: TH
Hai mặt cánh có độ chênh áp $500\,\text{Pa}$ trên diện tích $0,6\,\text{m}$ ^2. Coi lực nâng $F=\Delta pA$. Tính lực nâng.
\choiceTF
{\True Kết quả đúng là $300\,\text{N}$.}
{Kết quả bằng $600\,\text{N}$.}
{\True Khi hợp lực bằng 0, vật có gia tốc bằng 0.}
{Có thể bỏ qua đơn vị của đại lượng trong kết quả vật lí.}
\loigiai{
\begin{itemchoice}
\item (Đúng) Kết quả đúng là $300\,\text{N}$. (Vì): $F=\Delta pA=500\cdot0,6=300$ N. b) Sai vì giá trị hoặc chiều nêu ra không phù hợp. c) Đúng theo định luật II Newton. d) Sai vì kết quả vật lí phải kèm đơn vị thích hợp
\item (Sai) Kết quả bằng $600\,\text{N}$.
\item (Đúng) Khi hợp lực bằng 0, vật có gia tốc bằng 0.
\item (Sai) Có thể bỏ qua đơn vị của đại lượng trong kết quả vật lí.
\end{itemchoice}
}
\end{ex}

% ===== Câu 123 =====
% ID: AIQ_L10C3_FULL_1095
% Mức: VD
\begin{ex}
% ID: AIQ_L10C3_FULL_1095
% Mức: VD
Hai mặt cánh có độ chênh áp $550\,\text{Pa}$ trên diện tích $0,7\,\text{m}$ ^2. Coi lực nâng $F=\Delta pA$. Tính lực nâng.
\choiceTF
{\True Kết quả đúng là $385\,\text{N}$.}
{Kết quả bằng $770\,\text{N}$.}
{\True Khi hợp lực bằng 0, vật có gia tốc bằng 0.}
{Có thể bỏ qua đơn vị của đại lượng trong kết quả vật lí.}
\loigiai{
\begin{itemchoice}
\item (Đúng) Kết quả đúng là $385\,\text{N}$. (Vì): $F=\Delta pA=550\cdot0,7=385$ N. b) Sai vì giá trị hoặc chiều nêu ra không phù hợp. c) Đúng theo định luật II Newton. d) Sai vì kết quả vật lí phải kèm đơn vị thích hợp
\item (Sai) Kết quả bằng $770\,\text{N}$.
\item (Đúng) Khi hợp lực bằng 0, vật có gia tốc bằng 0.
\item (Sai) Có thể bỏ qua đơn vị của đại lượng trong kết quả vật lí.
\end{itemchoice}
}
\end{ex}

% ===== Câu 124 =====
% ID: AIQ_L10C3_FULL_1096
% Mức: TH
\begin{ex}
% ID: AIQ_L10C3_FULL_1096
% Mức: TH
Hai mặt cánh có độ chênh áp $200\,\text{Pa}$ trên diện tích $0,8\,\text{m}$ ^2. Coi lực nâng $F=\Delta pA$. Tính lực nâng.
\shortans{160}
\loigiai{
$F=\Delta pA=200\cdot0,8=160$ N. Đáp án trả lời ngắn không ghi đơn vị.
}
\end{ex}

% ===== Câu 125 =====
% ID: AIQ_L10C3_FULL_1097
% Mức: TH
\begin{ex}
% ID: AIQ_L10C3_FULL_1097
% Mức: TH
Hai mặt cánh có độ chênh áp $250\,\text{Pa}$ trên diện tích $0,9\,\text{m}$ ^2. Coi lực nâng $F=\Delta pA$. Tính lực nâng.
\shortans{225}
\loigiai{
$F=\Delta pA=250\cdot0,9=225$ N. Đáp án trả lời ngắn không ghi đơn vị.
}
\end{ex}

% ===== Câu 126 =====
% ID: AIQ_L10C3_FULL_1098
% Mức: VD
\begin{ex}
% ID: AIQ_L10C3_FULL_1098
% Mức: VD
Hai mặt cánh có độ chênh áp $300\,\text{Pa}$ trên diện tích $0,4\,\text{m}$ ^2. Coi lực nâng $F=\Delta pA$. Tính lực nâng.
\shortans{120}
\loigiai{
$F=\Delta pA=300\cdot0,4=120$ N. Đáp án trả lời ngắn không ghi đơn vị.
}
\end{ex}

% ===== Câu 127 =====
% ID: AIQ_L10C3_FULL_1099
% Mức: VD
\begin{ex}
% ID: AIQ_L10C3_FULL_1099
% Mức: VD
Hai mặt cánh có độ chênh áp $350\,\text{Pa}$ trên diện tích $0,5\,\text{m}$ ^2. Coi lực nâng $F=\Delta pA$. Tính lực nâng.
\shortans{175}
\loigiai{
$F=\Delta pA=350\cdot0,5=175$ N. Đáp án trả lời ngắn không ghi đơn vị.
}
\end{ex}

% ===== Câu 128 =====
% ID: AIQ_L10C3_FULL_1100
% Mức: VD
\begin{ex}
% ID: AIQ_L10C3_FULL_1100
% Mức: VD
Hai mặt cánh có độ chênh áp $400\,\text{Pa}$ trên diện tích $0,6\,\text{m}$ ^2. Coi lực nâng $F=\Delta pA$. Tính lực nâng.
\shortans{240}
\loigiai{
$F=\Delta pA=400\cdot0,6=240$ N. Đáp án trả lời ngắn không ghi đơn vị.
}
\end{ex}

% ===== Câu 129 =====
% ID: AIQ_L10C3_FULL_1101
% Mức: VDC
\begin{ex}
% ID: AIQ_L10C3_FULL_1101
% Mức: VDC
Hai mặt cánh có độ chênh áp $450\,\text{Pa}$ trên diện tích $0,7\,\text{m}$ ^2. Coi lực nâng $F=\Delta pA$. Tính lực nâng.
\shortans{315}
\loigiai{
$F=\Delta pA=450\cdot0,7=315$ N. Đáp án trả lời ngắn không ghi đơn vị.
}
\end{ex}

% ===== Câu 130 =====
% ID: AIQ_L10C3_FULL_1102
% Mức: TH
\begin{bt}
% ID: AIQ_L10C3_FULL_1102
% Mức: TH
Hai mặt cánh có độ chênh áp $500\,\text{Pa}$ trên diện tích $0,8\,\text{m}$ ^2. Coi lực nâng $F=\Delta pA$. Tính lực nâng.
\loigiai{
Phân tích dữ kiện, chọn định luật phù hợp. $F=\Delta pA=500\cdot0,8=400$ N.
}
\end{bt}

% ===== Câu 131 =====
% ID: AIQ_L10C3_FULL_1103
% Mức: TH
\begin{bt}
% ID: AIQ_L10C3_FULL_1103
% Mức: TH
Hai mặt cánh có độ chênh áp $550\,\text{Pa}$ trên diện tích $0,9\,\text{m}$ ^2. Coi lực nâng $F=\Delta pA$. Tính lực nâng.
\loigiai{
Phân tích dữ kiện, chọn định luật phù hợp. $F=\Delta pA=550\cdot0,9=495$ N.
}
\end{bt}

% ===== Câu 132 =====
% ID: AIQ_L10C3_FULL_1104
% Mức: VD
\begin{bt}
% ID: AIQ_L10C3_FULL_1104
% Mức: VD
Hai mặt cánh có độ chênh áp $200\,\text{Pa}$ trên diện tích $0,4\,\text{m}$ ^2. Coi lực nâng $F=\Delta pA$. Tính lực nâng.
\loigiai{
Phân tích dữ kiện, chọn định luật phù hợp. $F=\Delta pA=200\cdot0,4=80$ N.
}
\end{bt}

% ===== Câu 133 =====
% ID: AIQ_L10C3_FULL_1105
% Mức: VD
\begin{bt}
% ID: AIQ_L10C3_FULL_1105
% Mức: VD
Hai mặt cánh có độ chênh áp $250\,\text{Pa}$ trên diện tích $0,5\,\text{m}$ ^2. Coi lực nâng $F=\Delta pA$. Tính lực nâng.
\loigiai{
Phân tích dữ kiện, chọn định luật phù hợp. $F=\Delta pA=250\cdot0,5=125$ N.
}
\end{bt}

% ===== Câu 134 =====
% ID: AIQ_L10C3_FULL_1106
% Mức: VD
\begin{bt}
% ID: AIQ_L10C3_FULL_1106
% Mức: VD
Hai mặt cánh có độ chênh áp $300\,\text{Pa}$ trên diện tích $0,6\,\text{m}$ ^2. Coi lực nâng $F=\Delta pA$. Tính lực nâng.
\loigiai{
Phân tích dữ kiện, chọn định luật phù hợp. $F=\Delta pA=300\cdot0,6=180$ N.
}
\end{bt}

% ===== Câu 135 =====
% ID: AIQ_L10C3_FULL_1107
% Mức: VDC
\begin{bt}
% ID: AIQ_L10C3_FULL_1107
% Mức: VDC
Hai mặt cánh có độ chênh áp $350\,\text{Pa}$ trên diện tích $0,7\,\text{m}$ ^2. Coi lực nâng $F=\Delta pA$. Tính lực nâng.
\loigiai{
Phân tích dữ kiện, chọn định luật phù hợp. $F=\Delta pA=350\cdot0,7=245$ N.
}
\end{bt}

% ===== Câu 136 =====
% ID: AIQ_L10C3_FULL_1108
% Mức: NB
\dangbt{Vận dụng lực cản và lực nâng trong thực tiễn}
\begin{ex}
% ID: AIQ_L10C3_FULL_1108
% Mức: NB
Lực cản tỉ lệ với bình phương tốc độ. Khi tốc độ tăng từ $12\,\text{m}$ /s lên $14\,\text{m}$ /s, lực cản tăng 1,36 lần. Kết quả này được suy ra thế nào?
\choice
{\True 1,36 lần}
{1,63 lần}
{1,09 lần}
{3,36 lần}
\loigiai{
$F_2/F_1=(v_2/v_1)^2=[(14)/12]^2=1,36$. Chọn A.
}
\end{ex}

% ===== Câu 137 =====
% ID: AIQ_L10C3_FULL_1109
% Mức: NB
\begin{ex}
% ID: AIQ_L10C3_FULL_1109
% Mức: NB
Lực cản tỉ lệ với bình phương tốc độ. Khi tốc độ tăng từ $13\,\text{m}$ /s lên $15\,\text{m}$ /s, lực cản tăng 1,33 lần. Kết quả này được suy ra thế nào?
\choice
{1,6 lần}
{1,06 lần}
{3,33 lần}
{\True 1,33 lần}
\loigiai{
$F_2/F_1=(v_2/v_1)^2=[(15)/13]^2=1,33$. Chọn D.
}
\end{ex}

% ===== Câu 138 =====
% ID: AIQ_L10C3_FULL_1110
% Mức: NB
\begin{ex}
% ID: AIQ_L10C3_FULL_1110
% Mức: NB
Lực cản tỉ lệ với bình phương tốc độ. Khi tốc độ tăng từ $14\,\text{m}$ /s lên $16\,\text{m}$ /s, lực cản tăng 1,31 lần. Kết quả này được suy ra thế nào?
\choice
{1,05 lần}
{3,31 lần}
{\True 1,31 lần}
{1,57 lần}
\loigiai{
$F_2/F_1=(v_2/v_1)^2=[(16)/14]^2=1,31$. Chọn C.
}
\end{ex}

% ===== Câu 139 =====
% ID: AIQ_L10C3_FULL_1111
% Mức: TH
\begin{ex}
% ID: AIQ_L10C3_FULL_1111
% Mức: TH
Lực cản tỉ lệ với bình phương tốc độ. Khi tốc độ tăng từ $15\,\text{m}$ /s lên $17\,\text{m}$ /s, lực cản tăng 1,28 lần. Kết quả này được suy ra thế nào?
\choice
{3,28 lần}
{\True 1,28 lần}
{1,54 lần}
{1,02 lần}
\loigiai{
$F_2/F_1=(v_2/v_1)^2=[(17)/15]^2=1,28$. Chọn B.
}
\end{ex}

% ===== Câu 140 =====
% ID: AIQ_L10C3_FULL_1112
% Mức: TH
\begin{ex}
% ID: AIQ_L10C3_FULL_1112
% Mức: TH
Lực cản tỉ lệ với bình phương tốc độ. Khi tốc độ tăng từ $16\,\text{m}$ /s lên $18\,\text{m}$ /s, lực cản tăng 1,27 lần. Kết quả này được suy ra thế nào?
\choice
{\True 1,27 lần}
{1,52 lần}
{1,02 lần}
{3,27 lần}
\loigiai{
$F_2/F_1=(v_2/v_1)^2=[(18)/16]^2=1,27$. Chọn A.
}
\end{ex}

% ===== Câu 141 =====
% ID: AIQ_L10C3_FULL_1113
% Mức: TH
\begin{ex}
% ID: AIQ_L10C3_FULL_1113
% Mức: TH
Lực cản tỉ lệ với bình phương tốc độ. Khi tốc độ tăng từ $10\,\text{m}$ /s lên $12\,\text{m}$ /s, lực cản tăng 1,44 lần. Kết quả này được suy ra thế nào?
\choice
{1,73 lần}
{1,15 lần}
{3,44 lần}
{\True 1,44 lần}
\loigiai{
$F_2/F_1=(v_2/v_1)^2=[(12)/10]^2=1,44$. Chọn D.
}
\end{ex}

% ===== Câu 142 =====
% ID: AIQ_L10C3_FULL_1114
% Mức: TH
\begin{ex}
% ID: AIQ_L10C3_FULL_1114
% Mức: TH
Lực cản tỉ lệ với bình phương tốc độ. Khi tốc độ tăng từ $11\,\text{m}$ /s lên $13\,\text{m}$ /s, lực cản tăng 1,4 lần. Kết quả này được suy ra thế nào?
\choice
{1,12 lần}
{3,4 lần}
{\True 1,4 lần}
{1,68 lần}
\loigiai{
$F_2/F_1=(v_2/v_1)^2=[(13)/11]^2=1,4$. Chọn C.
}
\end{ex}

% ===== Câu 143 =====
% ID: AIQ_L10C3_FULL_1115
% Mức: VD
\begin{ex}
% ID: AIQ_L10C3_FULL_1115
% Mức: VD
Lực cản tỉ lệ với bình phương tốc độ. Khi tốc độ tăng từ $12\,\text{m}$ /s lên $14\,\text{m}$ /s, lực cản tăng 1,36 lần. Kết quả này được suy ra thế nào?
\choice
{3,36 lần}
{\True 1,36 lần}
{1,63 lần}
{1,09 lần}
\loigiai{
$F_2/F_1=(v_2/v_1)^2=[(14)/12]^2=1,36$. Chọn B.
}
\end{ex}

% ===== Câu 144 =====
% ID: AIQ_L10C3_FULL_1116
% Mức: VD
\begin{ex}
% ID: AIQ_L10C3_FULL_1116
% Mức: VD
Lực cản tỉ lệ với bình phương tốc độ. Khi tốc độ tăng từ $13\,\text{m}$ /s lên $15\,\text{m}$ /s, lực cản tăng 1,33 lần. Kết quả này được suy ra thế nào?
\choice
{\True 1,33 lần}
{1,6 lần}
{1,06 lần}
{3,33 lần}
\loigiai{
$F_2/F_1=(v_2/v_1)^2=[(15)/13]^2=1,33$. Chọn A.
}
\end{ex}

% ===== Câu 145 =====
% ID: AIQ_L10C3_FULL_1117
% Mức: VDC
\begin{ex}
% ID: AIQ_L10C3_FULL_1117
% Mức: VDC
Lực cản tỉ lệ với bình phương tốc độ. Khi tốc độ tăng từ $14\,\text{m}$ /s lên $16\,\text{m}$ /s, lực cản tăng 1,31 lần. Kết quả này được suy ra thế nào?
\choice
{1,57 lần}
{1,05 lần}
{3,31 lần}
{\True 1,31 lần}
\loigiai{
$F_2/F_1=(v_2/v_1)^2=[(16)/14]^2=1,31$. Chọn D.
}
\end{ex}

% ===== Câu 146 =====
% ID: AIQ_L10C3_FULL_1118
% Mức: NB
\begin{ex}
% ID: AIQ_L10C3_FULL_1118
% Mức: NB
Lực cản tỉ lệ với bình phương tốc độ. Khi tốc độ tăng từ $15\,\text{m}$ /s lên $17\,\text{m}$ /s, lực cản tăng 1,28 lần. Kết quả này được suy ra thế nào?
\choiceTF
{\True Kết quả đúng là 1,28 lần.}
{Kết quả bằng 2,56 lần.}
{\True Khi hợp lực bằng 0, vật có gia tốc bằng 0.}
{Có thể bỏ qua đơn vị của đại lượng trong kết quả vật lí.}
\loigiai{
\begin{itemchoice}
\item (Đúng) Kết quả đúng là 1,28 lần. (Vì): $F_2/F_1=(v_2/v_1)^2=[(17)/15]^2=1,28$. b) Sai vì giá trị hoặc chiều nêu ra không phù hợp. c) Đúng theo định luật II Newton. d) Sai vì kết quả vật lí phải kèm đơn vị thích hợp
\item (Sai) Kết quả bằng 2,56 lần.
\item (Đúng) Khi hợp lực bằng 0, vật có gia tốc bằng 0.
\item (Sai) Có thể bỏ qua đơn vị của đại lượng trong kết quả vật lí.
\end{itemchoice}
}
\end{ex}

% ===== Câu 147 =====
% ID: AIQ_L10C3_FULL_1119
% Mức: NB
\begin{ex}
% ID: AIQ_L10C3_FULL_1119
% Mức: NB
Lực cản tỉ lệ với bình phương tốc độ. Khi tốc độ tăng từ $16\,\text{m}$ /s lên $18\,\text{m}$ /s, lực cản tăng 1,27 lần. Kết quả này được suy ra thế nào?
\choiceTF
{\True Kết quả đúng là 1,27 lần.}
{Kết quả bằng 2,54 lần.}
{\True Khi hợp lực bằng 0, vật có gia tốc bằng 0.}
{Có thể bỏ qua đơn vị của đại lượng trong kết quả vật lí.}
\loigiai{
\begin{itemchoice}
\item (Đúng) Kết quả đúng là 1,27 lần. (Vì): $F_2/F_1=(v_2/v_1)^2=[(18)/16]^2=1,27$. b) Sai vì giá trị hoặc chiều nêu ra không phù hợp. c) Đúng theo định luật II Newton. d) Sai vì kết quả vật lí phải kèm đơn vị thích hợp
\item (Sai) Kết quả bằng 2,54 lần.
\item (Đúng) Khi hợp lực bằng 0, vật có gia tốc bằng 0.
\item (Sai) Có thể bỏ qua đơn vị của đại lượng trong kết quả vật lí.
\end{itemchoice}
}
\end{ex}

% ===== Câu 148 =====
% ID: AIQ_L10C3_FULL_1120
% Mức: TH
\begin{ex}
% ID: AIQ_L10C3_FULL_1120
% Mức: TH
Lực cản tỉ lệ với bình phương tốc độ. Khi tốc độ tăng từ $10\,\text{m}$ /s lên $12\,\text{m}$ /s, lực cản tăng 1,44 lần. Kết quả này được suy ra thế nào?
\choiceTF
{\True Kết quả đúng là 1,44 lần.}
{Kết quả bằng 2,88 lần.}
{\True Khi hợp lực bằng 0, vật có gia tốc bằng 0.}
{Có thể bỏ qua đơn vị của đại lượng trong kết quả vật lí.}
\loigiai{
\begin{itemchoice}
\item (Đúng) Kết quả đúng là 1,44 lần. (Vì): $F_2/F_1=(v_2/v_1)^2=[(12)/10]^2=1,44$. b) Sai vì giá trị hoặc chiều nêu ra không phù hợp. c) Đúng theo định luật II Newton. d) Sai vì kết quả vật lí phải kèm đơn vị thích hợp
\item (Sai) Kết quả bằng 2,88 lần.
\item (Đúng) Khi hợp lực bằng 0, vật có gia tốc bằng 0.
\item (Sai) Có thể bỏ qua đơn vị của đại lượng trong kết quả vật lí.
\end{itemchoice}
}
\end{ex}

% ===== Câu 149 =====
% ID: AIQ_L10C3_FULL_1121
% Mức: TH
\begin{ex}
% ID: AIQ_L10C3_FULL_1121
% Mức: TH
Lực cản tỉ lệ với bình phương tốc độ. Khi tốc độ tăng từ $11\,\text{m}$ /s lên $13\,\text{m}$ /s, lực cản tăng 1,4 lần. Kết quả này được suy ra thế nào?
\choiceTF
{\True Kết quả đúng là 1,4 lần.}
{Kết quả bằng 2,8 lần.}
{\True Khi hợp lực bằng 0, vật có gia tốc bằng 0.}
{Có thể bỏ qua đơn vị của đại lượng trong kết quả vật lí.}
\loigiai{
\begin{itemchoice}
\item (Đúng) Kết quả đúng là 1,4 lần. (Vì): $F_2/F_1=(v_2/v_1)^2=[(13)/11]^2=1,4$. b) Sai vì giá trị hoặc chiều nêu ra không phù hợp. c) Đúng theo định luật II Newton. d) Sai vì kết quả vật lí phải kèm đơn vị thích hợp
\item (Sai) Kết quả bằng 2,8 lần.
\item (Đúng) Khi hợp lực bằng 0, vật có gia tốc bằng 0.
\item (Sai) Có thể bỏ qua đơn vị của đại lượng trong kết quả vật lí.
\end{itemchoice}
}
\end{ex}

% ===== Câu 150 =====
% ID: AIQ_L10C3_FULL_1122
% Mức: VD
\begin{ex}
% ID: AIQ_L10C3_FULL_1122
% Mức: VD
Lực cản tỉ lệ với bình phương tốc độ. Khi tốc độ tăng từ $12\,\text{m}$ /s lên $14\,\text{m}$ /s, lực cản tăng 1,36 lần. Kết quả này được suy ra thế nào?
\choiceTF
{\True Kết quả đúng là 1,36 lần.}
{Kết quả bằng 2,72 lần.}
{\True Khi hợp lực bằng 0, vật có gia tốc bằng 0.}
{Có thể bỏ qua đơn vị của đại lượng trong kết quả vật lí.}
\loigiai{
\begin{itemchoice}
\item (Đúng) Kết quả đúng là 1,36 lần. (Vì): $F_2/F_1=(v_2/v_1)^2=[(14)/12]^2=1,36$. b) Sai vì giá trị hoặc chiều nêu ra không phù hợp. c) Đúng theo định luật II Newton. d) Sai vì kết quả vật lí phải kèm đơn vị thích hợp
\item (Sai) Kết quả bằng 2,72 lần.
\item (Đúng) Khi hợp lực bằng 0, vật có gia tốc bằng 0.
\item (Sai) Có thể bỏ qua đơn vị của đại lượng trong kết quả vật lí.
\end{itemchoice}
}
\end{ex}

% ===== Câu 151 =====
% ID: AIQ_L10C3_FULL_1123
% Mức: TH
\begin{ex}
% ID: AIQ_L10C3_FULL_1123
% Mức: TH
Lực cản tỉ lệ với bình phương tốc độ. Khi tốc độ tăng từ $13\,\text{m}$ /s lên $15\,\text{m}$ /s, lực cản tăng 1,33 lần. Kết quả này được suy ra thế nào?
\shortans{1,33}
\loigiai{
$F_2/F_1=(v_2/v_1)^2=[(15)/13]^2=1,33$. Đáp án trả lời ngắn không ghi đơn vị.
}
\end{ex}

% ===== Câu 152 =====
% ID: AIQ_L10C3_FULL_1124
% Mức: TH
\begin{ex}
% ID: AIQ_L10C3_FULL_1124
% Mức: TH
Lực cản tỉ lệ với bình phương tốc độ. Khi tốc độ tăng từ $14\,\text{m}$ /s lên $16\,\text{m}$ /s, lực cản tăng 1,31 lần. Kết quả này được suy ra thế nào?
\shortans{1,31}
\loigiai{
$F_2/F_1=(v_2/v_1)^2=[(16)/14]^2=1,31$. Đáp án trả lời ngắn không ghi đơn vị.
}
\end{ex}

% ===== Câu 153 =====
% ID: AIQ_L10C3_FULL_1125
% Mức: VD
\begin{ex}
% ID: AIQ_L10C3_FULL_1125
% Mức: VD
Lực cản tỉ lệ với bình phương tốc độ. Khi tốc độ tăng từ $15\,\text{m}$ /s lên $17\,\text{m}$ /s, lực cản tăng 1,28 lần. Kết quả này được suy ra thế nào?
\shortans{1,28}
\loigiai{
$F_2/F_1=(v_2/v_1)^2=[(17)/15]^2=1,28$. Đáp án trả lời ngắn không ghi đơn vị.
}
\end{ex}

% ===== Câu 154 =====
% ID: AIQ_L10C3_FULL_1126
% Mức: VD
\begin{ex}
% ID: AIQ_L10C3_FULL_1126
% Mức: VD
Lực cản tỉ lệ với bình phương tốc độ. Khi tốc độ tăng từ $16\,\text{m}$ /s lên $18\,\text{m}$ /s, lực cản tăng 1,27 lần. Kết quả này được suy ra thế nào?
\shortans{1,27}
\loigiai{
$F_2/F_1=(v_2/v_1)^2=[(18)/16]^2=1,27$. Đáp án trả lời ngắn không ghi đơn vị.
}
\end{ex}

% ===== Câu 155 =====
% ID: AIQ_L10C3_FULL_1127
% Mức: VD
\begin{ex}
% ID: AIQ_L10C3_FULL_1127
% Mức: VD
Lực cản tỉ lệ với bình phương tốc độ. Khi tốc độ tăng từ $10\,\text{m}$ /s lên $12\,\text{m}$ /s, lực cản tăng 1,44 lần. Kết quả này được suy ra thế nào?
\shortans{1,44}
\loigiai{
$F_2/F_1=(v_2/v_1)^2=[(12)/10]^2=1,44$. Đáp án trả lời ngắn không ghi đơn vị.
}
\end{ex}

% ===== Câu 156 =====
% ID: AIQ_L10C3_FULL_1128
% Mức: VDC
\begin{ex}
% ID: AIQ_L10C3_FULL_1128
% Mức: VDC
Lực cản tỉ lệ với bình phương tốc độ. Khi tốc độ tăng từ $11\,\text{m}$ /s lên $13\,\text{m}$ /s, lực cản tăng 1,4 lần. Kết quả này được suy ra thế nào?
\shortans{1,4}
\loigiai{
$F_2/F_1=(v_2/v_1)^2=[(13)/11]^2=1,4$. Đáp án trả lời ngắn không ghi đơn vị.
}
\end{ex}

% ===== Câu 157 =====
% ID: AIQ_L10C3_FULL_1129
% Mức: TH
\begin{bt}
% ID: AIQ_L10C3_FULL_1129
% Mức: TH
Lực cản tỉ lệ với bình phương tốc độ. Khi tốc độ tăng từ $12\,\text{m}$ /s lên $14\,\text{m}$ /s, lực cản tăng 1,36 lần. Kết quả này được suy ra thế nào?
\loigiai{
Phân tích dữ kiện, chọn định luật phù hợp. $F_2/F_1=(v_2/v_1)^2=[(14)/12]^2=1,36$.
}
\end{bt}

% ===== Câu 158 =====
% ID: AIQ_L10C3_FULL_1130
% Mức: TH
\begin{bt}
% ID: AIQ_L10C3_FULL_1130
% Mức: TH
Lực cản tỉ lệ với bình phương tốc độ. Khi tốc độ tăng từ $13\,\text{m}$ /s lên $15\,\text{m}$ /s, lực cản tăng 1,33 lần. Kết quả này được suy ra thế nào?
\loigiai{
Phân tích dữ kiện, chọn định luật phù hợp. $F_2/F_1=(v_2/v_1)^2=[(15)/13]^2=1,33$.
}
\end{bt}

% ===== Câu 159 =====
% ID: AIQ_L10C3_FULL_1131
% Mức: VD
\begin{bt}
% ID: AIQ_L10C3_FULL_1131
% Mức: VD
Lực cản tỉ lệ với bình phương tốc độ. Khi tốc độ tăng từ $14\,\text{m}$ /s lên $16\,\text{m}$ /s, lực cản tăng 1,31 lần. Kết quả này được suy ra thế nào?
\loigiai{
Phân tích dữ kiện, chọn định luật phù hợp. $F_2/F_1=(v_2/v_1)^2=[(16)/14]^2=1,31$.
}
\end{bt}

% ===== Câu 160 =====
% ID: AIQ_L10C3_FULL_1132
% Mức: VD
\begin{bt}
% ID: AIQ_L10C3_FULL_1132
% Mức: VD
Lực cản tỉ lệ với bình phương tốc độ. Khi tốc độ tăng từ $15\,\text{m}$ /s lên $17\,\text{m}$ /s, lực cản tăng 1,28 lần. Kết quả này được suy ra thế nào?
\loigiai{
Phân tích dữ kiện, chọn định luật phù hợp. $F_2/F_1=(v_2/v_1)^2=[(17)/15]^2=1,28$.
}
\end{bt}

% ===== Câu 161 =====
% ID: AIQ_L10C3_FULL_1133
% Mức: VD
\begin{bt}
% ID: AIQ_L10C3_FULL_1133
% Mức: VD
Lực cản tỉ lệ với bình phương tốc độ. Khi tốc độ tăng từ $16\,\text{m}$ /s lên $18\,\text{m}$ /s, lực cản tăng 1,27 lần. Kết quả này được suy ra thế nào?
\loigiai{
Phân tích dữ kiện, chọn định luật phù hợp. $F_2/F_1=(v_2/v_1)^2=[(18)/16]^2=1,27$.
}
\end{bt}

% ===== Câu 162 =====
% ID: AIQ_L10C3_FULL_1134
% Mức: VDC
\begin{bt}
% ID: AIQ_L10C3_FULL_1134
% Mức: VDC
Lực cản tỉ lệ với bình phương tốc độ. Khi tốc độ tăng từ $10\,\text{m}$ /s lên $12\,\text{m}$ /s, lực cản tăng 1,44 lần. Kết quả này được suy ra thế nào?
\loigiai{
Phân tích dữ kiện, chọn định luật phù hợp. $F_2/F_1=(v_2/v_1)^2=[(12)/10]^2=1,44$.
}
\end{bt}
